\section{プログラムソース}

\subsection{報告者担当分}
報告者が担当した，トロンボーンおよびカスタネットの音源のソースコードを示す．いずれもProcessingで記述されており，クライアントからシリアル通信で受信した音符の情報をもとに音を合成して再生する．また，合成音の波形を描画し，基本周波数を測定して表示する機能を備える．

\subsubsection{トロンボーンの音源}

\begin{lstlisting}
import processing.serial.*;
import ddf.minim.*;
import ddf.minim.ugens.*;

Serial myPort;
Minim minim;
AudioOutput out;
Waveform currentWaveform;
// --- 表示用変数 ---
String currentNote = "REST";
int currentBpm = -1; // 初期値は未設定にして外部から読み込む
String currentDurationName = "REST";

// 基準音符長（デフォルト: 1.0 = 1拍単位）
float ToneLength = 1.0f;

float[] maxAmp = new float[29];
long lastNoteEndMillis = 0;
float currentToneLengthSec = 0.0f; // 現在鳴っている音の秒数
float currentToneLengthBeats = 0.0f; // 現在鳴っている音の拍数

// --- 周波数計測用 ---
float measuredFreq = 0.0f;   // 自己相関法で計測した基本周波数(Hz)
float displayedFreq = 0.0f;  // 表示用に平滑化した自己相関周波数(Hz)
float measuredZc = 0.0f;     // ゼロ交差法で計測した基本周波数(Hz)
float displayedZc = 0.0f;    // 表示用に平滑化したゼロ交差周波数(Hz)
float displayedPeak = 0.0f;  // 表示用に平滑化したピーク振幅
float AMP_GATE = 0.01f;      // この音量(RMS)未満は無音とみなし計測しない

// 目標周波数：C4(261.63Hz)を1オクターブ下げた値（本システムはfreq/2で発音）
float TARGET_FREQ = 261.63f / 2.0f; // = 130.815 Hz（C3相当）

// --- トロンボーンの特性を反映したクラス ---
class HackInstrument implements Instrument
{
  Oscil wave;
  MoogFilter filter; 
  Delay delay;
  ADSR ampEnv;
  Line freqEnv;
  float baseFreq;
  float maxAmpValue;

  HackInstrument(float frequency, float amplitude, Waveform wf)
  {
    baseFreq = frequency;
    maxAmpValue = amplitude;
    wave = new Oscil(frequency * 0.96f, 1.0, wf);

    filter = new MoogFilter(1200.0f, 0.2f, MoogFilter.Type.LP);
    delay = new Delay(0.2f, 0.3f, true, true);

    // ADSRの引数: (最大音量, アタック秒, ディケイ秒, サステイン比率, リリース秒)
    // 管楽器らしく 0.08秒かけて音が立ち上がり、吹いている間は 80% の音量をキープします
    ampEnv = new ADSR(amplitude, 0.08f, 0.05f, 0.8f, 0.3f);
    
    // 波形の出力を音量エンベロープ（ampEnv）にパッチする
    wave.patch(filter).patch(delay).patch(ampEnv);
    
    freqEnv = new Line();
    freqEnv.patch(wave.frequency);
  }

  void noteOn(float duration)
  {
    ampEnv.noteOn(); // ADSRのトリガー開始
    ampEnv.patch(out); // 出力ポートへ接続
    freqEnv.activate(0.12f, baseFreq * 0.96f, baseFreq);
  }

  void noteOff()
  {
    ampEnv.noteOff(); // 音を消し始める（リリースタイム開始）
    ampEnv.unpatchAfterRelease(out); // 音が完全に消え去った後に自動で接続を解除（ブツ切り防止）
  }
}

void setup()
{
  size(600, 300);
  pixelDensity(1); // 高密度ディスプレイでの自動2倍描画を無効化（従来の見た目に戻す）
  // 日本語が「□（豆腐）」にならないよう、日本語対応フォントを設定する
  textFont(pickJapaneseFont(40));
  textSize(24);
  minim = new Minim(this);
  out = minim.getLineOut();
  out.setGain(2.0);
  currentBpm = loadBpm();
  out.setTempo(currentBpm);
  
  // シリアルポートの初期化処理
  // 利用可能なポートを表示し、Arduino（usbmodem / usbserial）を自動で探して接続する
  printArray(Serial.list());
  String portName = null;
  for (int i = 0; i < Serial.list().length; i++) {
    String s = Serial.list()[i];
    if (s.indexOf("usbmodem") >= 0 || s.indexOf("usbserial") >= 0) {
      portName = s;
      break;
    }
  }
  if (portName != null) {
    myPort = new Serial(this, portName, 115200);
    myPort.bufferUntil('\n');
    println("Serial connected: " + portName);
  } else {
    println("Arduinoのポートが見つかりません。上のリストから手動で指定してください。");
  }
  
  for(int i=0; i<maxAmp.length; i++) maxAmp[i] = 0.15f;
  
  setTromboneWave();
  textSize(24);
}

void setTromboneWave() {
  currentWaveform = WavetableGenerator.gen10(
    4096, 
    new float[] {
      1.0000f, 0.96f, 0.90f, 0.84f, 0.76f, 
      0.72f, 0.71f, 0.63f, 0.57f, 0.56f
    }
  );
}

// BPMを外部から取得する: 環境変数 HCK_BPM または BPM、
// もしくはスケッチの data フォルダに置いた "bpm.txt" の1行目を読み込む。
int loadBpm() {
  // 1) 環境変数
  try {
    String env = System.getenv("HCK_BPM");
    if (env == null) env = System.getenv("BPM");
    if (env != null) {
      env = trim(env);
      if (env.length() > 0) {
        try { return Integer.parseInt(env); } catch (Exception e) {}
      }
    }
  } catch (Exception e) {
    // ignore
  }

  // 2) data/bpm.txt
  try {
    String[] lines = loadStrings("bpm.txt");
    if (lines != null && lines.length > 0) {
      String s = trim(lines[0]);
      if (s.length() > 0) {
        try { return Integer.parseInt(s); } catch (Exception e) {}
      }
    }
  } catch (Exception e) {
    // ファイルが無ければ無視
  }

  // 3) フォールバック
  return 120;
}

void draw() {
  background(0); // 画面を黒で塗りつぶしてリセット（オシロスコープ風）

  // --- オシロスコープ風グリッド（目盛り）と軸の単位を描画 ---
  // 縦方向の拡大率：振幅1.0が gain ピクセルに対応する
  float gain = height * 1.5f;
  // 画面横幅ぶんに相当する時間(ms)＝バッファ長 / サンプリングレート
  float totalMs = out.bufferSize() / out.sampleRate() * 1000.0f;
  drawScopeGrid(gain, totalMs);

  // 出力バッファ（out.mix）の波形を描画（正方向を上＝オシロスコープと同じ向き）
  stroke(0, 255, 120); // 波形の色（緑）
  strokeWeight(1.5f);
  noFill();
  for (int i = 0; i < out.bufferSize() - 1; i++) {
    float x1 = map(i,     0, out.bufferSize(), 0, width);
    float x2 = map(i + 1, 0, out.bufferSize(), 0, width);
    float y1 = height / 2 - out.mix.get(i)     * gain;
    float y2 = height / 2 - out.mix.get(i + 1) * gain;
    line(x1, y1, x2, y2);
  }

  // --- 周波数を計測して画面に表示 ---
  measuredFreq = detectFrequency(out.mix, out.sampleRate());        // 自己相関法
  measuredZc   = detectFreqZeroCross(out.mix, out.sampleRate());    // ゼロ交差法
  float peak   = peakAmplitude(out.mix);                            // ピーク振幅
  if (measuredFreq > 0) {
    // 急な変動を抑えるため指数移動平均で平滑化
    displayedFreq = (displayedFreq <= 0) ? measuredFreq
                                         : displayedFreq * 0.8f + measuredFreq * 0.2f;
    displayedZc   = (displayedZc <= 0 || measuredZc <= 0) ? measuredZc
                                         : displayedZc * 0.8f + measuredZc * 0.2f;
    displayedPeak = displayedPeak * 0.8f + peak * 0.2f;
  } else {
    displayedFreq = 0;
    displayedZc = 0;
    displayedPeak = 0; // 無音
  }
  drawFreqReadout(displayedFreq, displayedZc, displayedPeak);

  // 保存リクエストがあれば、ヒント等を描く前のきれいな画面を保存する
  if (saveRequested) {
    String fn = "capture_" + timestamp() + ".png";
    save(fn); // スケッチフォルダ直下に保存（波形＋周波数のみ）
    lastSavedName = fn;
    savedFlashUntil = millis() + 1500;
    println("画像を保存しました: " + sketchPath(fn));
    saveRequested = false;
  }
}

// オシロスコープ風：薄い目盛りグリッド＋原点を通る太い縦軸/横軸を描画する
void drawScopeGrid(float gain, float totalMs) {
  pushStyle();
  textSize(11);

  float cy = height / 2.0f;      // 振幅0（横軸）のy
  float ox = 2;                  // 時間0（縦軸）のx（左端。太線が見えるよう少し内側）
  float msPerDiv = 5.0f;
  float ampPerDiv = 0.1f;
  float ampMax = cy / gain;      // 画面端に相当する振幅
  int   maxDiv  = (int)(ampMax / ampPerDiv);

  // ===== 1) 薄い目盛りグリッド（背景・控えめ）=====
  strokeWeight(1);
  stroke(38);
  for (float t = 0; t <= totalMs + 0.001f; t += msPerDiv) {   // 時間の縦線
    float x = t / totalMs * width;
    line(x, 0, x, height);
  }
  for (int k = -maxDiv; k <= maxDiv; k++) {                   // 振幅の横線
    if (k == 0) continue;
    float y = cy - k * ampPerDiv * gain;
    line(0, y, width, y);
  }

  // ===== 2) 原点を通る太い軸（はっきり）=====
  stroke(0, 210, 255);   // 水色で目立たせる（波形の緑と区別）
  strokeWeight(2);
  line(0, cy, width, cy);      // 横軸（時間 t：振幅0のライン）
  line(ox, 0, ox, height);     // 縦軸（振幅 A：時間0のライン）

  // 軸上の目盛りマーク（ティック）
  stroke(0, 210, 255);
  strokeWeight(1);
  for (float t = msPerDiv; t < totalMs; t += msPerDiv) {      // 横軸のティック
    float x = t / totalMs * width;
    line(x, cy - 4, x, cy + 4);
  }
  for (int k = -maxDiv; k <= maxDiv; k++) {                   // 縦軸のティック
    if (k == 0) continue;
    float y = cy - k * ampPerDiv * gain;
    line(ox, y, ox + 8, y);
  }

  // ===== 3) 数値ラベル =====
  noStroke();
  textAlign(CENTER, TOP);
  fill(190);
  for (float t = msPerDiv; t < totalMs; t += msPerDiv) {      // 時間[ms]（横軸の下）
    float x = t / totalMs * width;
    text(nf(t, 0, 0), x, cy + 6);
  }
  textAlign(RIGHT, CENTER);
  for (int k = -maxDiv; k <= maxDiv; k++) {                   // 振幅（右端）
    if (k == 0) continue;
    float a = k * ampPerDiv;
    float y = cy - a * gain;
    text((a > 0 ? "+" : "") + nf(a, 0, 1), width - 6, y);
  }

  // ===== 4) 軸タイトル・原点 =====
  fill(0, 210, 255);
  textAlign(RIGHT, TOP);
  text("時間 t [ms] →", width - 6, cy + 22);        // 横軸タイトル
  textAlign(LEFT, TOP);
  text("↑ 振幅 A（正規化, ±1.0）", ox + 6, 110);     // 縦軸タイトル（左上の周波数パネルの下に置く）
  textAlign(LEFT, TOP);
  fill(190);
  text("O", ox + 4, cy + 4);                        // 原点

  // 目盛り幅の凡例（左下）
  textAlign(LEFT, BOTTOM);
  text("1目盛 = 横 " + nf(msPerDiv, 0, 0) + " ms ／ 縦 " + nf(ampPerDiv, 0, 1)
       + "（フルスケール±1.0）", 6, height - 2);

  popStyle();
}

// 画面左上に各計測値（自己相関・ゼロ交差・ピーク振幅・目標誤差）を表示する
void drawFreqReadout(float fAuto, float fZc, float peak) {
  pushStyle();
  textAlign(LEFT, TOP);

  // 文字が波形と重なっても読めるよう、背面に半透明の暗いパネルを敷く
  noStroke();
  fill(0, 170); // 黒・不透明度170/255（波形がうっすら透ける程度）
  rect(0, 0, 300, (fAuto > 0) ? 104 : 50);

  if (fAuto > 0) {
    // 自己相関法（メイン）
    fill(0, 255, 120);
    textSize(26);
    text(nf(fAuto, 0, 2) + " Hz", 10, 5);

    String name = nearestNoteName(fAuto);
    float cents = centsOff(fAuto);
    fill(200);
    textSize(13);
    text("自己相関  " + name + " (" + sign(cents) + nf(cents, 0, 1) + " cent)", 10, 36);

    // ゼロ交差法
    if (fZc > 0) {
      text("ゼロ交差  " + nf(fZc, 0, 2) + " Hz", 10, 52);
    } else {
      text("ゼロ交差  -- Hz", 10, 52);
    }

    // ピーク振幅（1.0超でクリップ）
    text("ピーク振幅  " + nf(peak, 0, 3) + (peak >= 1.0f ? "（クリップ）" : "（クリップなし）"), 10, 68);

    // 目標値に対する誤差（自己相関法ベース）
    float errHz = fAuto - TARGET_FREQ;
    float errPct = errHz / TARGET_FREQ * 100.0f;
    fill(errPct >= -1.0f && errPct <= 1.0f ? color(0, 255, 120) : color(255, 120, 120));
    text("目標 " + nf(TARGET_FREQ, 0, 2) + "Hz  誤差 "
         + sign(errHz) + nf(errHz, 0, 2) + "Hz / " + sign(errPct) + nf(errPct, 0, 2) + "%",
         10, 84);
  } else {
    fill(120);
    textSize(26);
    text("-- Hz", 10, 5);
    textSize(13);
    text("（無音）", 10, 36);
  }
  popStyle();
}

// 符号付き表示用（+を明示）
String sign(float v) {
  return v >= 0 ? "+" : "";
}

// 自己相関による基本周波数の推定（低音域でも精度が出る）
float detectFrequency(AudioBuffer buf, float sampleRate) {
  int n = buf.size();

  // 1) 音量(RMS)を見て、無音なら計測しない
  float rms = 0;
  for (int i = 0; i < n; i++) {
    float v = buf.get(i);
    rms += v * v;
  }
  rms = sqrt(rms / n);
  if (rms < AMP_GATE) return 0;

  // 2) 探索する周期(ラグ)の範囲：約 50Hz 〜 1500Hz
  int minLag = (int)(sampleRate / 1500.0f);
  int maxLag = (int)(sampleRate / 50.0f);
  if (maxLag > n - 1) maxLag = n - 1;
  if (minLag < 2) minLag = 2;

  // 3) 自己相関が最大になるラグを探す
  float bestCorr = 0;
  int bestLag = -1;
  for (int lag = minLag; lag <= maxLag; lag++) {
    float corr = 0;
    for (int i = 0; i < n - lag; i++) {
      corr += buf.get(i) * buf.get(i + lag);
    }
    if (corr > bestCorr) {
      bestCorr = corr;
      bestLag = lag;
    }
  }
  if (bestLag <= 0) return 0;

  // 4) 放物線補間でサブサンプル精度に補正
  float lag = bestLag;
  if (bestLag > minLag && bestLag < maxLag) {
    float c0 = autocorrAt(buf, bestLag - 1);
    float c1 = autocorrAt(buf, bestLag);
    float c2 = autocorrAt(buf, bestLag + 1);
    float denom = (c0 - 2 * c1 + c2);
    if (abs(denom) > 1e-9) {
      lag = bestLag + 0.5f * (c0 - c2) / denom;
    }
  }
  return sampleRate / lag;
}

// 指定ラグでの自己相関値（補間用ヘルパー）
float autocorrAt(AudioBuffer buf, int lag) {
  int n = buf.size();
  float corr = 0;
  for (int i = 0; i < n - lag; i++) {
    corr += buf.get(i) * buf.get(i + lag);
  }
  return corr;
}

// ゼロ交差法による基本周波数の推定
// 倍音のさざ波を誤検出しないよう、ピーク振幅の±25%のヒステリシス（シュミットトリガ）で
// 立ち上がりエッジのみを数え、ゼロ点を線形補間して最初と最後の交差から平均周期を求める
float detectFreqZeroCross(AudioBuffer buf, float sampleRate) {
  int n = buf.size();
  float peak = peakAmplitude(buf);
  if (peak < AMP_GATE) return 0;

  float hi = peak * 0.25f;   // 上側しきい値
  float lo = -peak * 0.25f;  // 下側しきい値
  boolean wasLow = false;    // 直近で下側しきい値を下回ったか
  float firstCross = -1, lastCross = -1;
  int cycles = 0;

  for (int i = 1; i < n; i++) {
    float v = buf.get(i);
    if (v < lo) wasLow = true;
    if (wasLow && v > hi) {
      // 立ち上がり確定。直前サンプルとの間でゼロ点(=0)を線形補間
      float v0 = buf.get(i - 1);
      float pos = i;
      if (v != v0) pos = (i - 1) + (0 - v0) / (v - v0);
      if (firstCross < 0) {
        firstCross = pos;
      } else {
        lastCross = pos;
        cycles++;
      }
      wasLow = false;
    }
  }
  if (cycles < 1 || lastCross < 0) return 0;
  float periodSamples = (lastCross - firstCross) / cycles;
  if (periodSamples <= 0) return 0;
  return sampleRate / periodSamples;
}

// バッファ内の最大振幅（絶対値）
float peakAmplitude(AudioBuffer buf) {
  int n = buf.size();
  float peak = 0;
  for (int i = 0; i < n; i++) {
    float a = abs(buf.get(i));
    if (a > peak) peak = a;
  }
  return peak;
}

// 周波数に最も近い音名（A4=440の12平均律）
String nearestNoteName(float f) {
  String[] names = {"C", "C#", "D", "D#", "E", "F", "F#", "G", "G#", "A", "A#", "B"};
  int midi = round(69 + 12 * (log(f / 440.0f) / log(2)));
  int octave = midi / 12 - 1;
  return names[(midi % 12 + 12) % 12] + octave;
}

// 最寄り音からのズレ（セント, ±50が境目）
float centsOff(float f) {
  float midi = 69 + 12 * (log(f / 440.0f) / log(2));
  float nearest = round(midi);
  return (midi - nearest) * 100.0f;
}

// --- 画面（波形＋周波数表示）を画像として保存 ---
// S キー: 現在の画面を PNG でスケッチフォルダに保存する
int savedFlashUntil = 0;        // 「保存しました」表示を消す時刻(millis)
String lastSavedName = "";
boolean saveRequested = false;  // S キーで立てる保存フラグ（実際の保存は draw 内で行う）

void keyPressed() {
  if (key == 's' || key == 'S') {
    saveRequested = true; // 描画タイミングに合わせて draw() 内で保存する
  }
}

// 日本語を表示できるフォントを、この環境にあるものから自動で選ぶ
// （見つからなければ既定フォントにフォールバック）
PFont pickJapaneseFont(int sz) {
  String[] avail = PFont.list();
  // Mac/Win/Linux の代表的な日本語フォント名のキーワード（部分一致で探す）
  String[] keys = {
    "Hiragino", "YuGothic", "Yu Gothic", "Noto Sans CJK", "Noto Sans JP",
    "Noto Serif CJK", "MS Gothic", "ＭＳ", "Meiryo", "メイリオ", "Osaka", "IPAGothic", "TakaoGothic"
  };
  for (String k : keys) {
    for (String a : avail) {
      if (a.toLowerCase().indexOf(k.toLowerCase()) >= 0) {
        println("日本語フォントを使用: " + a);
        return createFont(a, sz, true);
      }
    }
  }
  println("日本語フォントが見つかりませんでした。既定フォントを使用します。");
  return createFont(avail.length > 0 ? avail[0] : "SansSerif", sz, true);
}

// ファイル名用のタイムスタンプ（例: 20260630_071530）
String timestamp() {
  return year()
       + nf(month(), 2) + nf(day(), 2) + "_"
       + nf(hour(), 2) + nf(minute(), 2) + nf(second(), 2);
}

void serialEvent(Serial p) {
  String data = p.readStringUntil('\n');
  if (data != null) {
    data = trim(data);
    String[] values = split(data, ',');

    // クライアントの送信形式：「周波数,鳴らすミリ秒,強さ,BPM」 例: 262,500,100,120
    // "RAW_DATA:5" や "【BPM同期】…" などのデバッグ行はカンマ区切りで4つに分かれないので自動的に無視される
    if (values.length >= 4) {
      println(data); // 受信ノート記録（正常発音数テスト用の受信ログ。不要なら削除）
      try {
        // 1. 周波数(Hz)・ミリ秒・強さ・BPMを、それぞれ数値として読み込む
        float freq = float(values[0]) / 2.0f; // 周波数を半分にして1オクターブ下げる
        float durationMs = float(values[1]);
        // values[2] = velocity（強さ）。今回は音量に反映しないため未使用
        int bpmValue = int(float(values[3]));

        // BPMをクライアントに同期（表示と拍数計算用）
        if (bpmValue > 0) {
          currentBpm = bpmValue;
          out.setTempo(currentBpm);
        }

        // 2. 周波数が0（休符：REST）より大きいときだけ音を鳴らす
        if (freq > 0) {
          // 3. ミリ秒を1000で割って「秒」に変換する
          float durationSec = durationMs / 1000.0f;

          currentNote = freq + " Hz";
          currentToneLengthSec = durationSec;
          currentToneLengthBeats = durationMs / (60000.0f / currentBpm);
          lastNoteEndMillis = millis() + (long)durationMs;

        // 4. HackInstrumentに直接周波数（freq）と秒数を渡す
        out.playNote(0.0f, durationSec, new HackInstrument(freq, 0.15f, currentWaveform));
        }
      }
    catch (Exception e) {
      // 例外が発生した場合は無視して続行
    }
    }
  }
}
\end{lstlisting}

\subsubsection{カスタネットの音源}

\begin{lstlisting}
import processing.serial.*;
import ddf.minim.*;
import ddf.minim.ugens.*;

Serial myPort;
Minim minim;
AudioOutput out;
Waveform currentWaveform;

// --- 表示用変数 ---
String currentNote = "REST";
int currentBpm = -1; // 初期値は未設定にして外部から読み込む
String currentDurationName = "REST";

// 基準音符長（デフォルト: 1.0 = 1拍単位）
float ToneLength = 1.0f;

float[] maxAmp = new float[29];
long lastNoteEndMillis = 0;
float currentToneLengthSec = 0.0f;   // 現在鳴っている音の秒数
float currentToneLengthBeats = 0.0f; // 現在鳴っている音の拍数

// --- カスタネットの特性を反映したクラス ---
// --- カスタネット：Arduino連動 ＆ ユーザー微調整反映版 ---
class WoodPercussion implements Instrument
{
  // 1. 表面の硬い音（アタック成分）
  Noise clickNoise;
  ADSR clickEnv;
  MoogFilter clickFilter; 

  // 2. 木の内部の響き（ボディ成分）
  Noise bodyNoise;
  ADSR bodyEnv;
  MoogFilter bodyFilter;

  // 3. 高域の乾いた成分（理想に近い参考版に合わせて追加）
  Noise dryNoise;
  ADSR dryEnv;
  MoogFilter dryFilter;

  WoodPercussion(float inputFreq, float amplitude)
  {
    // 参考版の音作り：周波数を木の共鳴帯域に寄せ、毎回わずかに揺らす
    float bodyFreq = constrain(inputFreq, 1000.0f, 2600.0f);
    bodyFreq += random(-90.0f, 90.0f);
    float clickFreq = bodyFreq + random(2500.0f, 3600.0f);
    float dryFreq = bodyFreq + random(4200.0f, 5200.0f);
    float ampRand = amplitude * random(0.80f, 1.05f);

    // ①【表面の硬い音】カチッという入りを残す
    clickNoise = new Noise(ampRand * 1.15f, Noise.Tint.WHITE);
    clickEnv = new ADSR(1.0f, 0.0005f, 0.007f, 0.0f, 0.003f);
    clickFilter = new MoogFilter(clickFreq, 0.65f, MoogFilter.Type.BP);
    clickNoise.patch(clickEnv).patch(clickFilter);

    // ②【木の胴鳴り成分】「カッ」だけでなく「コッ」と響く感じを足す
    bodyNoise = new Noise(ampRand * 0.40f, Noise.Tint.PINK);
    bodyEnv = new ADSR(1.0f, 0.001f, 0.060f, 0.0f, 0.012f);
    bodyFilter = new MoogFilter(bodyFreq, 0.85f, MoogFilter.Type.BP);
    bodyNoise.patch(bodyEnv).patch(bodyFilter);

    // ③【高域のパキッとした成分】硬さを残しつつ余韻を作る
    dryNoise = new Noise(ampRand * 0.33f, Noise.Tint.WHITE);
    dryEnv = new ADSR(1.0f, 0.0005f, 0.018f, 0.0f, 0.006f);
    dryFilter = new MoogFilter(dryFreq, 0.55f, MoogFilter.Type.BP);
    dryNoise.patch(dryEnv).patch(dryFilter);
  }

  void noteOn(float duration)
  {
    clickEnv.noteOn();
    bodyEnv.noteOn();
    dryEnv.noteOn();
    clickFilter.patch(out);
    bodyFilter.patch(out);
    dryFilter.patch(out);
  }

  void noteOff()
  {
    clickEnv.noteOff();
    bodyEnv.noteOff();
    dryEnv.noteOff();
    clickFilter.unpatch(out);
    bodyFilter.unpatch(out);
    dryFilter.unpatch(out);
  }
}

void setup()
{
  size(600, 300);
  pixelDensity(1); // 高密度ディスプレイでの自動2倍描画を無効化（従来の見た目に戻す）
  // 日本語が「□（豆腐）」に文字化けしないよう、日本語対応フォントを設定する
  textFont(pickJapaneseFont(40));
  textSize(24);
  minim = new Minim(this);
  out = minim.getLineOut();
  out.setTempo(currentBpm);
  
  // シリアルポートの初期化処理
  // 利用可能なポートを表示し、Arduino（usbmodem / usbserial）を自動で探して接続する
  printArray(Serial.list());
  String portName = null;
  for (int i = 0; i < Serial.list().length; i++) {
    String s = Serial.list()[i];
    if (s.indexOf("usbmodem") >= 0 || s.indexOf("usbserial") >= 0) {
      portName = s;
      break;
    }
  }
  if (portName != null) {
    myPort = new Serial(this, portName, 115200);
    myPort.bufferUntil('\n');
    println("Serial connected: " + portName);
  } else {
    println("Arduinoのポートが見つかりません。上のリストから手動で指定してください。");
  }
  
  for(int i=0; i<maxAmp.length; i++) maxAmp[i] = 0.15f;
  
  setCastanetWave();
  textSize(24);
}

// カスタネットの「カチッ」とした乾いたクラック音を再現するため、高次倍音を豊富に含んだ硬い波形を生成
void setCastanetWave() {
  currentWaveform = WavetableGenerator.gen10(
    4096, 
    new float[] { 0.1f, 0.3f, 0.5f, 0.8f, 1.0f, 0.9f, 0.8f, 0.7f, 0.6f, 0.5f }
  );
}

void draw() {
  background(0); // 画面を黒で塗りつぶしてリセット（オシロスコープ風）

  // --- オシロスコープ風グリッド（目盛り）と軸の単位を描画 ---
  float gain = height * 1.5f;                                    // 振幅1.0が gain ピクセル
  float totalMs = out.bufferSize() / out.sampleRate() * 1000.0f; // 画面横幅ぶんの時間(ms)
  drawScopeGrid(gain, totalMs);

  // 出力バッファ（out.mix）の波形を描画（正方向を上＝オシロスコープと同じ向き）
  stroke(0, 255, 120); // 波形の色（緑）
  strokeWeight(1.5f);
  noFill();
  for (int i = 0; i < out.bufferSize() - 1; i++) {
    float x1 = map(i,     0, out.bufferSize(), 0, width);
    float x2 = map(i + 1, 0, out.bufferSize(), 0, width);
    float y1 = height / 2 - out.mix.get(i)     * gain;
    float y2 = height / 2 - out.mix.get(i + 1) * gain;
    line(x1, y1, x2, y2);
  }
}

// 日本語を表示できるフォントを、この環境にあるものから自動で選ぶ
PFont pickJapaneseFont(int sz) {
  String[] avail = PFont.list();
  String[] keys = {
    "Hiragino", "YuGothic", "Yu Gothic", "Noto Sans CJK", "Noto Sans JP",
    "Noto Serif CJK", "MS Gothic", "Meiryo", "メイリオ", "Osaka", "IPAGothic", "TakaoGothic"
  };
  for (String k : keys) {
    for (String a : avail) {
      if (a.toLowerCase().indexOf(k.toLowerCase()) >= 0) {
        println("日本語フォントを使用: " + a);
        return createFont(a, sz, true);
      }
    }
  }
  println("日本語フォントが見つかりませんでした。既定フォントを使用します。");
  return createFont(avail.length > 0 ? avail[0] : "SansSerif", sz, true);
}

// オシロスコープ風：薄い目盛りグリッド＋原点を通る太い縦軸/横軸を描画する
void drawScopeGrid(float gain, float totalMs) {
  pushStyle();
  textSize(11);

  float cy = height / 2.0f;      // 振幅0（横軸）のy
  float ox = 2;                  // 時間0（縦軸）のx（左端。太線が見えるよう少し内側）
  float msPerDiv = 5.0f;
  float ampPerDiv = 0.1f;
  float ampMax = cy / gain;      // 画面端に相当する振幅
  int   maxDiv  = (int)(ampMax / ampPerDiv);

  // ===== 1) 薄い目盛りグリッド（背景・控えめ）=====
  strokeWeight(1);
  stroke(38);
  for (float t = 0; t <= totalMs + 0.001f; t += msPerDiv) {   // 時間の縦線
    float x = t / totalMs * width;
    line(x, 0, x, height);
  }
  for (int k = -maxDiv; k <= maxDiv; k++) {                   // 振幅の横線
    if (k == 0) continue;
    float y = cy - k * ampPerDiv * gain;
    line(0, y, width, y);
  }

  // ===== 2) 原点を通る太い軸（はっきり）=====
  stroke(0, 210, 255);   // 水色で目立たせる（波形の緑と区別）
  strokeWeight(2);
  line(0, cy, width, cy);      // 横軸（時間 t：振幅0のライン）
  line(ox, 0, ox, height);     // 縦軸（振幅 A：時間0のライン）

  // 軸上の目盛りマーク（ティック）
  stroke(0, 210, 255);
  strokeWeight(1);
  for (float t = msPerDiv; t < totalMs; t += msPerDiv) {      // 横軸のティック
    float x = t / totalMs * width;
    line(x, cy - 4, x, cy + 4);
  }
  for (int k = -maxDiv; k <= maxDiv; k++) {                   // 縦軸のティック
    if (k == 0) continue;
    float y = cy - k * ampPerDiv * gain;
    line(ox, y, ox + 8, y);
  }

  // ===== 3) 数値ラベル =====
  noStroke();
  textAlign(CENTER, TOP);
  fill(190);
  for (float t = msPerDiv; t < totalMs; t += msPerDiv) {      // 時間[ms]（横軸の下）
    float x = t / totalMs * width;
    text(nf(t, 0, 0), x, cy + 6);
  }
  textAlign(RIGHT, CENTER);
  for (int k = -maxDiv; k <= maxDiv; k++) {                   // 振幅（右端）
    if (k == 0) continue;
    float a = k * ampPerDiv;
    float y = cy - a * gain;
    text((a > 0 ? "+" : "") + nf(a, 0, 1), width - 6, y);
  }

  // ===== 4) 軸タイトル・原点 =====
  fill(0, 210, 255);
  textAlign(RIGHT, TOP);
  text("時間 t [ms] →", width - 6, cy + 22);        // 横軸タイトル
  textAlign(LEFT, TOP);
  text("↑ 振幅 A（正規化, ±1.0）", ox + 6, 8);       // 縦軸タイトル
  fill(190);
  text("O", ox + 4, cy + 4);                        // 原点

  // 目盛り幅の凡例（左下）
  textAlign(LEFT, BOTTOM);
  text("1目盛 = 横 " + nf(msPerDiv, 0, 0) + " ms ／ 縦 " + nf(ampPerDiv, 0, 1)
       + "（フルスケール±1.0）", 6, height - 2);

  popStyle();
}

void keyPressed() {
  // 画像保存（S）: 現在の波形画面を PNG で保存
  if (key == 's' || key == 'S') {
    String ts = year() + nf(month(), 2) + nf(day(), 2) + "_"
              + nf(hour(), 2) + nf(minute(), 2) + nf(second(), 2);
    save("castanet_capture_" + ts + ".png");
    println("画像を保存しました: castanet_capture_" + ts + ".png");
  }
}

void serialEvent(Serial p) {
  String data = p.readStringUntil('\n');
  if (data != null) {
    data = trim(data);
    String[] values = split(data, ',');

    // クライアントの送信形式：「周波数,鳴らすミリ秒,強さ,BPM」 例: 262,500,100,120
    // "RAW_DATA:5" や "【BPM同期】…" などのデバッグ行はカンマ区切りで4つに分かれないので自動的に無視される
    if (values.length >= 4) {
      println(data); // 受信ノート記録（正常発音数テスト用の受信ログ。不要なら削除）
      try {
        // 1. 周波数(Hz)・ミリ秒・強さ・BPMをそれぞれ数値としてパース
        // カスタネットらしい抜けの良い高域を引き出すため、送られてきた周波数をそのまま使用
        float freq = float(values[0]);
        float durationMs = float(values[1]);
        // values[2] = velocity（強さ）。今回は音量に反映しないため未使用
        int bpmValue = int(float(values[3]));

        // BPMをクライアントに同期（表示と拍数計算用）
        if (bpmValue > 0) {
          currentBpm = bpmValue;
          out.setTempo(currentBpm);
        }

        // 2. 周波数が0（休符）より大きいときだけ打音を生成
        if (freq > 0) {
          float durationSec = durationMs / 1000.0f;

          currentNote = freq + " Hz";
          currentToneLengthSec = durationSec;
          currentToneLengthBeats = durationMs / (60000.0f / currentBpm);
          lastNoteEndMillis = millis() + (long)durationMs;

          // 3. カスタネット特性を持たせたインストゥルメントオブジェクトを生成して再生
          out.playNote(0.0f, 0.8f, new WoodPercussion(freq, 1.0f));
        }
      }
      catch (Exception e) {
        // データ破損などの例外発生時は安全に処理をスキップ
      }
    }
  }
}
\end{lstlisting}

\subsection{担当箇所以外}
システムの根幹となる，サーバおよびクライアントのソースコードを示す．これらは主に他のメンバーが実装したものである．

\subsubsection{サーバ（server.ino）}

\begin{lstlisting}
#include <WiFiS3.h>
#include <WiFiUdp.h>
#include "server_function.h"
#include "Arduino_LED_Matrix.h" 

// 【ネットワーク設定】
char ssid[] = "hackathon003-WPA2";
char pass[] = "hackathon003";
IPAddress BCaddress(255, 255, 255, 255); // 全員に一斉送信（ブロードキャスト）

uint16_t Port = 3000;
WiFiUDP udp;

// 【システム設定用変数】
ArduinoLEDMatrix matrix; 

// --- 小節管理用 ---
uint32_t LastBarTime = 0;
uint16_t Interval = 2000;
uint8_t BarCount = 39;

// --- BPM管理用 ---
uint8_t BPM = 120;
uint32_t LastPressTime = 0;
bool Flag = false;


void setup() {
  Serial.begin(115200);
  matrix.begin();
  
  // スイッチピンの設定
  pinMode(2, INPUT); 
  pinMode(3, INPUT);

  // 初期Intervalの計算
  Interval = (60000 / BPM) * 4;
  
  // Wi-Fi接続処理
  Serial.print("サーバー：Wi-Fiに接続中: ");
  Serial.println(ssid);
  WiFi.begin(ssid, pass);
  while (WiFi.status() != WL_CONNECTED) {
    delay(500);
    Serial.print(".");
  }
  Serial.println("");

  // IPアドレス取得待機
  while (WiFi.localIP() == IPAddress(0, 0, 0, 0)) {
    Serial.println("ルーターからIPアドレスを取得中...");
    delay(500);
  }
  
  Serial.println("接続完了!");
  Serial.print("サーバー自身のIPアドレス: ");
  Serial.println(WiFi.localIP());

  udp.begin(Port);
}


void loop() {
  Bar_control(&LastBarTime, Interval, &BarCount, BCaddress, Port, udp);
  BPM_control(&BPM, &Interval, &LastPressTime, &Flag, BCaddress, Port, udp);
}
\end{lstlisting}
サーバのメインプログラムである．\texttt{setup()}でWi-Fiへの接続とスイッチピンの初期化を行い，\texttt{loop()}で小節番号の送信（\texttt{Bar\_control}）とテンポの変更・送信（\texttt{BPM\_control}）を繰り返し実行する．
\subsubsection{サーバの関数ヘッダ（server\_function.h）}

\begin{lstlisting}
#ifndef SERVER_FUNCTION_H
#define SERVER_FUNCTION_H
#include <Arduino.h>
#include <WiFiUdp.h>
#include "Arduino_LED_Matrix.h"

void Bar_control(uint32_t *LastBarTime, uint16_t Interval, uint8_t *BarCount, IPAddress BCaddress, uint16_t Port, WiFiUDP &udp);

void BPM_control(uint8_t *BPM, uint16_t *Interval, uint32_t *LastPressTime, bool *Flag, IPAddress BCaddress, uint16_t Port, WiFiUDP &udp);

#endif
\end{lstlisting}
サーバの主要な関数のプロトタイプ宣言である．
\subsubsection{サーバの関数実装（server\_function.cpp）}

\begin{lstlisting}
#include "server_function.h"

// ============================================================
//  スイッチ押下検出ピン
// ============================================================
#define PIN_BPM_UP    2   // BPMアップ
#define PIN_BPM_DOWN  3   // BPMダウン
// ============================================================
//  デバウンスタイム
// ============================================================
#define DEBOUNCE_MS 30

void Bar_control(uint32_t *LastBarTime, uint16_t Interval, uint8_t *BarCount, IPAddress BCaddress, uint16_t Port, WiFiUDP &udp) {
  uint32_t temp = millis();
  
  // Interval（間隔）以上時間が経過したかチェック
  if (temp - *LastBarTime >= Interval) {
    *LastBarTime = temp;
    
    // 小節番号のカウント（上限39）
    if (*BarCount >= 39) {
      *BarCount = 0;
    } else {
      *BarCount += 1;
    }
    
    // UDPブロードキャスト送信
    udp.beginPacket(BCaddress, Port);
    udp.write(*BarCount);
    udp.endPacket();

    // デバッグ用出力
    Serial.print("送信した小節番号: ");
    Serial.println(*BarCount);
  }
}
// ------------------------------------------------------------
//  チャタリング制御用関数
// ------------------------------------------------------------
static bool detectPress(uint8_t pin, int *lastReading, int *stableState, uint32_t *lastChangeTime) {
  int reading = digitalRead(pin);

  // 読み値が変化したら、変化時刻を更新
  if (reading != *lastReading) {
    *lastChangeTime = millis();
    *lastReading = reading;
  }

  // 一定時間読み値が安定していたら確定状態を更新する
  if (millis() - *lastChangeTime > DEBOUNCE_MS) {
    if (reading != *stableState) {
      *stableState = reading;
      // 確定状態が HIGH に変わった瞬間 ＝ 押された瞬間
      if (*stableState == HIGH) {
        return true;
      }
    }
  }
  return false;
}


void BPM_control(uint8_t *BPM, uint16_t *Interval, uint32_t *LastPressTime, bool *Flag, IPAddress BCaddress, uint16_t Port, WiFiUDP &udp) {

  // --- 各ボタンのデバウンス状態（関数をまたいで保持するため static） ---
  static int      upLastReading   = LOW;
  static int      upStableState   = LOW;
  static uint32_t upLastChange    = 0;

  static int      downLastReading = LOW;
  static int      downStableState = LOW;
  static uint32_t downLastChange  = 0;

  // === BPMアップ（D2）===
  if (detectPress(PIN_BPM_UP, &upLastReading, &upStableState, &upLastChange)) {
    if (*BPM <= 250) {              // 上限255を超えない範囲で（+5してもオーバーフローしない）
      *BPM += 5;
      *Flag = true;
      *LastPressTime = millis();    // 最後に押した時刻
      Serial.print("BPM UP! 現在のBPM: "); Serial.println(*BPM);
    }
  }

  // === BPMダウン（D3）===
  if (detectPress(PIN_BPM_DOWN, &downLastReading, &downStableState, &downLastChange)) {
    if (*BPM >= 45) {               // 下限40を下回らない範囲で（-5してもアンダーフローしない）
      *BPM -= 5;
      *Flag = true;
      *LastPressTime = millis();
      Serial.print("BPM DOWN! 現在のBPM: "); Serial.println(*BPM);
    }
  }

  // === BPM変更後の送信処理 ===
  // 最後の押下から1秒以上 何も押されなければ、まとめて1回だけ送信する
  if (*Flag) {
    if (millis() - *LastPressTime > 1000) {
      udp.beginPacket(BCaddress, Port);
      udp.write(*BPM);
      udp.endPacket();
      
      *Flag = false;
      
      // 新しいBPMから小節の時間間隔(Interval)を再計算
      // 1分=60000ミリ秒。4/4拍子なので4倍
      *Interval = (60000 / *BPM) * 4; 
      
      Serial.print("【送信完了】BPM値を全楽器に送信しました。新Interval: ");
      Serial.println(*Interval);
    }
  }
}

\end{lstlisting}
サーバの主要な処理を実装したファイルである．\texttt{Bar\_control}関数は，一定間隔ごとに小節番号を0から39の範囲で循環させながらUDPで送信する．\texttt{detectPress}関数は，スイッチの読み値が30\,ms安定して初めて押下を確定することで，チャタリングによる誤検出を防ぐ．\texttt{BPM\_control}関数は，スイッチ押下に応じてテンポを5ずつ増減させ，最後の押下から1秒経過後に新しいテンポを送信する．

\subsubsection{クライアント（client.ino）}
楽器クライアントのメインプログラムである．\texttt{setup()}でWi-Fiへの接続と譜面データの初期化を行い，\texttt{loop()}で受信データの判別（\texttt{Parse\_data}），テンポの更新（\texttt{BPM\_update}），演奏処理（\texttt{Performance}），および人形の駆動（\texttt{robot.update}）を繰り返し実行する．各クライアントは，輪唱を実現するための固有のオフセット値（\texttt{Offset}）を持ち，受信した小節番号にこれを加えることで自身の演奏開始位置をずらす．これにより，同一の譜面を複数のクライアントが時間差で演奏し，輪唱として成立させている．

\begin{lstlisting}
#include <WiFiS3.h>
#include <WiFiUdp.h>
#include "client_function.h"
#include "EntertainmentRobot.h"


//Offset変数とsetupScore関数を楽器に合わせて下さい．その他の，client_function.cpp，client_function.hはいじらなくても大丈夫です．


// --- ネットワーク設定 ---
char ssid[] = "hackathon003-WPA2";     
char pass[] = "hackathon003"; 
uint16_t Port = 3000;    // サーバーと同じポート番号

WiFiUDP Udp;

EntertainmentRobot robot;

// --- クライアントサイドのグローバル変数 ---
char Offset = -2;         // 輪唱に必要なオフセット値（楽器ごとに設定，2番手だと-2）
uint8_t Data = 255;        // 受信データ
uint8_t CurrentBPM = 120;// 現在のBPM（初期値はサーバーに合わせて120とする）
uint8_t CurrentBar = 255;  // 現在演奏中の小節番号
uint8_t NoteIndex = 255; // 現在の小節の中で、何番目の音符を鳴らしているかのインデックス
bool Flag = false;       // 受信データが小節番号かBPMかを判定するフラグ
float ToneLength = 0.0;  // 基準音符（4分音符）の音の長さ（ミリ秒）
uint32_t StartTime = 0;  // 前回の音符を鳴らした時間（millis）
uint16_t Interval = 0;   // 実際に演奏する音の長さ（ミリ秒）

Pfm Score[25];           // 楽譜（40小節分の配列）


// --- 周波数定義 ---
#define NOTE_C4  262 // ド
#define NOTE_D4  294 // レ
#define NOTE_E4  330 // ミ
#define NOTE_F4  349 // ファ
#define NOTE_G4  392 // ソ
#define NOTE_A4  440 // ラ
#define REST     0   // 休符

#define NOTE_C5  523 // 高いド
#define NOTE_D5  587 // 高いレ
#define NOTE_E5  659 // 高いミ
#define NOTE_F5  698 // 高いファ
#define NOTE_G5  784 // 高いソ
#define NOTE_A5  880 // 高いラ

void setupScore() {
  // 第0小節〜第7小節（1周目）
  Score[0] = {{ {NOTE_C4, 1.0, 92}, {NOTE_D4, 1.0, 96}, {NOTE_E4, 1.0, 100}, {NOTE_F4, 1.0, 104} }, 4};
  Score[1] = {{ {NOTE_E4, 1.0, 100}, {NOTE_D4, 1.0, 96}, {NOTE_C4, 1.0, 90}, {REST, 1.0, 0} }, 4};
  Score[2] = {{ {NOTE_E4, 1.0, 96}, {NOTE_F4, 1.0, 100}, {NOTE_G4, 1.0, 106}, {NOTE_A4, 1.0, 112} }, 4};
  Score[3] = {{ {NOTE_G4, 1.0, 106}, {NOTE_F4, 1.0, 100}, {NOTE_E4, 1.0, 94}, {REST, 1.0, 0} }, 4};
  Score[4] = {{ {NOTE_C4, 1.0, 104}, {REST, 1.0, 0}, {NOTE_C4, 1.0, 104}, {REST, 1.0, 0} }, 4};
  Score[5] = {{ {NOTE_C4, 1.0, 98}, {REST, 1.0, 0}, {NOTE_C4, 1.0, 106}, {REST, 1.0, 0} }, 4};
  Score[6] = {{ {NOTE_C4, 0.5, 88}, {NOTE_C4, 0.5, 92}, {NOTE_D4, 0.5, 96}, {NOTE_D4, 0.5, 100}, {NOTE_E4, 0.5, 104}, {NOTE_E4, 0.5, 108}, {NOTE_F4, 0.5, 112}, {NOTE_F4, 0.5, 116} }, 8};
  Score[7] = {{ {NOTE_E4, 1.0, 108}, {NOTE_D4, 1.0, 100}, {NOTE_C4, 1.0, 92}, {REST, 1.0, 0} }, 4};

  // 第8小節（休符）
  Score[8] = {{ {REST, 4.0, 0} }, 1};
  Score[9] = {{ {REST, 4.0, 0} }, 1};
  Score[10] = {{ {REST, 4.0, 0} }, 1};
  Score[11] = {{ {REST, 4.0, 0} }, 1};
  Score[12] = {{ {REST, 4.0, 0} }, 1};

  // 第9小節〜第16小節（2周目：1オクターブ上）
  Score[13]  = {{ {NOTE_C5, 1.0, 92}, {NOTE_D5, 1.0, 96}, {NOTE_E5, 1.0, 100}, {NOTE_F5, 1.0, 104} }, 4};
  Score[14] = {{ {NOTE_E5, 1.0, 100}, {NOTE_D5, 1.0, 96}, {NOTE_C5, 1.0, 90}, {REST, 1.0, 0} }, 4};
  Score[15] = {{ {NOTE_E5, 1.0, 96}, {NOTE_F5, 1.0, 100}, {NOTE_G5, 1.0, 106}, {NOTE_A5, 1.0, 112} }, 4};
  Score[16] = {{ {NOTE_G5, 1.0, 106}, {NOTE_F5, 1.0, 100}, {NOTE_E5, 1.0, 94}, {REST, 1.0, 0} }, 4};
  Score[17] = {{ {NOTE_C5, 1.0, 104}, {REST, 1.0, 0}, {NOTE_C5, 1.0, 104}, {REST, 1.0, 0} }, 4};
  Score[18] = {{ {NOTE_C5, 1.0, 98}, {REST, 1.0, 0}, {NOTE_C5, 1.0, 106}, {REST, 1.0, 0} }, 4};
  Score[19] = {{ {NOTE_C5, 0.5, 88}, {NOTE_C5, 0.5, 92}, {NOTE_D5, 0.5, 96}, {NOTE_D5, 0.5, 100}, {NOTE_E5, 0.5, 104}, {NOTE_E5, 0.5, 108}, {NOTE_F5, 0.5, 112}, {NOTE_F5, 0.5, 116} }, 8};
  Score[20] = {{ {NOTE_E5, 1.0, 108}, {NOTE_D5, 1.0, 100}, {NOTE_C5, 1.0, 92}, {REST, 1.0, 0} }, 4};

  // 第17小節〜第18小節（休符）
  Score[21] = {{ {REST, 4.0, 0} }, 1};
  Score[22] = {{ {REST, 4.0, 0} }, 1};
  Score[23] = {{ {REST, 4.0, 0} }, 1};
  Score[24] = {{ {REST, 4.0, 0} }, 1};

  
}


void setup() {
  Serial.begin(115200);

  Serial.println("Wi-Fiに接続中...");
  WiFi.begin(ssid, pass);
  while (WiFi.status() != WL_CONNECTED) {
    delay(500);
    Serial.print(".");
  }
  Serial.println("\nWi-Fi接続完了!");

  // IPアドレス取得待機
  while (WiFi.localIP() == IPAddress(0, 0, 0, 0)) {
    Serial.println("ルーターからIPアドレスを取得中...");
    delay(500);
  }
  Serial.print("クライアントのIPアドレス: ");
  Serial.println(WiFi.localIP());
  
  Udp.begin(Port);

  setupScore();
  ToneLength = 60000.0 / CurrentBPM;
  Serial.println("クライアント：譜面データの読み込みが完了しました．");
  robot.begin(ssid, pass, Port);
}

void loop() {
  Flag = Parse_data(&Data, Offset, Udp);
  
  // 2. データがBPMだった場合の同期処理
  BPM_update(&CurrentBPM, Flag, Data, &ToneLength);

  // 3. 演奏位置制御とPCへのデータ送信（今回追加）
  Performance(Data, &CurrentBar, &NoteIndex, &StartTime, &Interval, ToneLength, CurrentBPM, Score);
  robot.update();
}

\end{lstlisting}

\subsubsection{クライアントの関数ヘッダ（client\_function.h）}
譜面データの構造体（\texttt{NoteData}，\texttt{Pfm}）と，クライアントの主要な関数のプロトタイプを定義したファイルである．\texttt{NoteData}は1つの音符の周波数・長さ・強さを，\texttt{Pfm}は1小節分の音符列を表す．

\begin{lstlisting}
#ifndef CLIENT_FUNCTION_H
#define CLIENT_FUNCTION_H
#include <Arduino.h>
#include <WiFiUdp.h>

// --- 譜面データの構造体定義 ---
// 1つの音符の情報を表す構造体
struct NoteData {
  int pitch;      // 音の高さ（周波数 Hz）
  float length;   // 音の長さ（基準となる4分音符を 1.0 とした相対値）
  uint8_t velocity; // 音の強さ（0〜127）
};

// 1つの小節（Pfm）の情報を表す構造体
struct Pfm {
  NoteData notes[8]; // 1小節に最大8つの音符が入ると仮定
  int noteCount;     // この小節に実際に含まれる音符の数
};

bool Parse_data(uint8_t *Data, char Offset, WiFiUDP &Udp);// 受信データがBPM値ならtrue（フラグON），小節番号ならfalse（フラグOFF）を返す関数

void BPM_update(uint8_t *CurrentBPM, bool Flag, uint8_t Data, float *ToneLength);

void Performance(uint8_t Data, uint8_t *CurrentBar, uint8_t *NoteIndex, uint32_t *StartTime, uint16_t *Interval, float ToneLength, uint8_t CurrentBPM, Pfm *Score);
#endif

\end{lstlisting}

\subsubsection{クライアントの関数実装（client\_function.cpp）}
クライアントの主要な処理を実装したファイルである．\texttt{Parse\_data}関数は，受信した1バイトのデータを値の大小によってテンポと小節番号に判別し，小節番号にはオフセットを加える．\texttt{BPM\_update}関数は，受信したテンポから基準となる4分音符の長さを算出する．\texttt{Performance}関数は，受信した小節番号に対応する音符を譜面データから順に読み出し，周波数・長さ・強さ・テンポをシリアル通信でProcessingへ送信する．

\begin{lstlisting}
#include "client_function.h"

bool Parse_data(uint8_t *Data, char Offset, WiFiUDP &Udp) {
  int packetSize = Udp.parsePacket();
  
  // 1. パケットを受信したかどうかの判定
  if (packetSize) {
    uint8_t received = Udp.read();
    Udp.flush();

    // === 【生存確認用デバッグ】受信した生の数値をそのままProcessingへ送る ===
    Serial.print("RAW_DATA:");
    Serial.println(received);
    // ===================================================================
    
    // 2. 受信データが40以上であればBPM値と判定（フラグONを返す）
    if (received >= 25) {
      *Data = received;
      return true; // フラグON
    }
    // 3. 40未満であれば小節番号と判定（オフセットを適用し，フラグOFFを返す）
    else {
      // サーバーの小節番号とオフセットを足す
      int myBar = received + Offset; 
      
      // マイナスになる場合（まだ自分の出番ではない待機期間）
      if (myBar < 0) {
        *Data = 255; // 待機状態を維持するための無効な値を入れる
      } else {
        *Data = myBar; // 自分の小節番号として適用
      }
      return false; // フラグOFF
    }

  }
  return false;
}

void BPM_update(uint8_t *CurrentBPM, bool Flag, uint8_t Data, float *ToneLength) {
  // フラグがOFF（小節番号）の場合は何もせず終了
  if (!Flag) {
    return;
  } 
  // フラグがON（BPM値）の場合の更新処理
  else {
    *CurrentBPM = Data;
    
    // 基準となる4分音符の長さをミリ秒単位で計算 (60000 / BPM)
    *ToneLength = 60000.0 / (*CurrentBPM);
    
    // 確認用のシリアル出力
    Serial.print("【BPM同期】新しいBPM: ");
    Serial.print(*CurrentBPM);
    Serial.print(" -> 4分音符の基準長さ(ms): ");
    Serial.println(*ToneLength);
  }
}

void Performance(uint8_t Data, uint8_t *CurrentBar, uint8_t *NoteIndex, uint32_t *StartTime, uint16_t *Interval, float ToneLength, uint8_t CurrentBPM, Pfm *Score) {
  
  // 1. 新しい小節番号を受信した時のリセット処理
  // （Dataが40未満で、かつ現在演奏中の小節と異なる場合）
  if (Data < 24 && Data != *CurrentBar) {
    *CurrentBar = Data;    // 現在の小節番号を更新
    *NoteIndex = 0;        // 読み込む音符を最初(0番目)に戻す
    *StartTime = millis(); // 演奏開始タイマーをリセット
    *Interval = 0;         // 即座に最初の音が鳴るようにIntervalを0にする
  }

  // フリーズ防止
  if (*CurrentBar >= 40) {
    return; 
  }

  // 2. 音を鳴らすタイミング（時間が経過したか）の判定
  if (millis() - *StartTime >= *Interval) {
    
    // その小節にある音符の数(noteCount)に達するまで処理を行う
    if (*NoteIndex < Score[*CurrentBar].noteCount) {
      
      *StartTime = millis(); // 次の音の基準時間を現在時刻に更新
      
      // 譜面配列から「音の高さ(pitch)」と「音の長さの割合(length)」を取得
      int pitch = Score[*CurrentBar].notes[*NoteIndex].pitch;
      float length = Score[*CurrentBar].notes[*NoteIndex].length;
      uint8_t velocity = Score[*CurrentBar].notes[*NoteIndex].velocity;
      
      // 次の音までの間隔（Interval）を計算（ミリ秒）
      *Interval = ToneLength * length;
      
      // PC（Processing）へシリアル通信で送信
      // 形式：「周波数,鳴らすミリ秒,強さ,BPM」 (例： 262,500,100,120)
      Serial.print(pitch);
      Serial.print(",");
      Serial.print(*Interval);
      Serial.print(",");
      Serial.print(velocity);
      Serial.print(",");
      Serial.println(CurrentBPM);
      
      *NoteIndex = *NoteIndex + 1; // 次の音符へ進む
    }
  }
}

\end{lstlisting}

\subsubsection{人形制御クラスのヘッダ（EntertainmentRobot.h）}
演奏に同期して人形のサーボモータを駆動する\texttt{EntertainmentRobot}クラスを定義したファイルである．楽器の種類（ピアノ・トロンボーン・バイオリン・カスタネット）に応じたサーボの角度や動作パターン，および楽譜に対応した動作テーブルを保持する．人形側にも輪唱のためのオフセット（\texttt{PART\_BAR\_OFFSET}）が設定されている．

\begin{lstlisting}
#ifndef ENTERTAINMENT_ROBOT_H
#define ENTERTAINMENT_ROBOT_H

#include <Arduino.h>
#include <WiFiS3.h>
#include <WiFiUdp.h>
#include <Servo.h>

// =====================
// 楽器選択定義
// =====================
#define PIANO      1
#define TROMBONE   2
#define VIOLIN     3
#define CASTANET   4

// 使う人形に合わせてここを変更
//#define INSTRUMENT PIANO
#define INSTRUMENT TROMBONE
//#define INSTRUMENT VIOLIN
//define INSTRUMENT CASTANET

// =====================
// 楽譜データ駆動のための構造体・列挙型
// =====================
enum ScoreEventType {
  SCORE_END  = 0,
  SCORE_PLAY = 1,
  SCORE_REST = 2
};

enum NoteLengthUnit {
  LENGTH_SIXTEENTH = 1,
  LENGTH_EIGHTH    = 2,
  LENGTH_QUARTER   = 4,
  LENGTH_HALF      = 8,
  LENGTH_WHOLE     = 16
};

struct ServoScoreEvent {
  ScoreEventType type;
  uint8_t lengthUnits;
};

class EntertainmentRobot {
public:
    EntertainmentRobot();
    void begin(const char* ssid, const char* pass, uint16_t port);
    void update();

private:
    // 通信関連
    WiFiUDP _udp;
    uint16_t _localPort;

    // サーボオブジェクト
    Servo _rightArmServo;
    Servo _leftArmServo;
    Servo _extraServo;

    // ピン配置定数
    static const int RIGHT_ARM_PIN = 9;
    static const int LEFT_ARM_PIN  = 11; // 新コードに合わせて10番に修正
    static const int EXTRA_PIN      = 5;

    // サーボ角度定数
    static const int ARM_REST_ANGLE = 90;
    static const int PIANO_RIGHT_PLAY_ANGLE = 70;
    static const int PIANO_LEFT_PLAY_ANGLE  = 110;

    static const int TROMBONE_FORWARD_ANGLE = 65;
    static const int TROMBONE_BACK_ANGLE    = 115;

    static const int VIOLIN_BOW_LEFT_ANGLE  = 65;
    static const int VIOLIN_BOW_RIGHT_ANGLE = 115;

    static const int CASTANET_PLAY_ANGLE    = 70;

    static const int HEAD_CENTER_ANGLE      = 90;
    static const int HEAD_LEFT_ANGLE        = 90; // 70度に修正
    static const int HEAD_RIGHT_ANGLE       = 180; // 110度に修正

    static const int MOUTH_CLOSE_ANGLE      = 90; // 90度に修正
    static const int MOUTH_OPEN_ANGLE       = 180; // 60度に修正

    // 全休符・オフセット設定
    static const uint8_t SERVER_BAR_COUNT = 25;
    static const int PART_BAR_OFFSET = -2; // 必要に応じて変更してください
    static const bool WHOLE_REST_BAR_TABLE[SERVER_BAR_COUNT];

    // 楽譜テーブルの定数
    static const uint8_t MAX_EVENTS_PER_BAR = 8;
    static const uint8_t BAR_LENGTH_UNITS = 16;
    static const uint8_t QUARTER_LENGTH_UNITS = 4;
    static const ServoScoreEvent SCORE_TABLE[SERVER_BAR_COUNT][MAX_EVENTS_PER_BAR];

    // 状態管理変数
    uint8_t _currentBPM;
    uint8_t _currentBar;
    unsigned long _beatInterval;
    
    // 楽譜イベント駆動用のタイマー変数
    unsigned long _nextEventTime;
    uint8_t _scoreEventIndex;
    uint8_t _usedUnitsInBar;

    bool _barActive;
    bool _armDirection;
    bool _motionMutedByWholeRest; // 全休符ミュートフラグ
    bool _scoreRestActive;        // 小節内の部分休符フラグ
    bool _pianoUseRightArm;       // ピアノの左右腕交互フラグ

    // サーボパルス（タイマー）管理
    bool _rightArmReturnActive;
    bool _leftArmReturnActive;
    bool _extraReturnActive;
    unsigned long _rightArmReturnTime;
    unsigned long _leftArmReturnTime;
    unsigned long _extraReturnTime;

    // BPM変更エフェクト管理
    bool _emotionActive;
    unsigned long _emotionEndTime;
    unsigned long _lastEmotionMoveTime;
    bool _emotionDirection;
    static const int EMOTION_BEATS = 4;

    // 内部メソッド
    void setupServos();
    void receiveUdpData();
    void updateBPM(uint8_t newBPM);
    void calculateBeatInterval();
    unsigned long noteLengthToMs(uint8_t lengthUnits);
    unsigned long calculatePulseTime(unsigned long eventDuration);
    uint8_t toScoreBar(uint8_t serverBarNumber);
    bool isWholeRestBar(uint8_t scoreBarNumber);
    void startBarMotion(uint8_t barNumber);
    void updateNormalMotion();
    void triggerNextScoreEvent();
    void triggerScoreMotion(unsigned long eventDuration);
    void finishBarMotion();
    void startEmotion();
    void updateEmotionMotion();
    
    void triggerRightArmPulse(int angle, unsigned long pulseTime);
    void triggerLeftArmPulse(int angle, unsigned long pulseTime);
    void triggerExtraPulse(int angle, unsigned long pulseTime);
    void updateServoReturn();
    void stopAllMotionForWholeRest();
    void stopMainMotionForRest();
    void cancelServoReturns();
    void resetMainServos();
    void resetEmotionServo();
    bool timeReached(unsigned long targetTime);
};

#endif 
\end{lstlisting}

\subsubsection{人形制御クラスの実装（EntertainmentRobot.cpp）}
\texttt{EntertainmentRobot}クラスの実装である．受信した小節番号とテンポに基づいてサーボを駆動し，演奏に同期した人形の動作を行う．また，テンポ変更時には，通常の演奏動作とは別の演出動作を行う．

\begin{lstlisting}
#include "EntertainmentRobot.h"

// ==========================================
// 全休符テーブルの定義 (false: 演奏, true: 全休符)
// ==========================================
const bool EntertainmentRobot::WHOLE_REST_BAR_TABLE[SERVER_BAR_COUNT] = {
  false, false, false, false, false, false, false, false, // 0〜7小節 (1周目)
  true,  true,  true,  true,  true,                       // 8〜12小節 (休符)
  false, false, false, false, false, false, false, false, // 13〜20小節 (2周目)
  true,  true,  true,  true                               // 21〜24小節 (休符)
};

// ==========================================
// 楽譜テーブルの定義 (25小節分)
// ==========================================
#define SCORE_FILL         { SCORE_END, 0 }

#define BAR_QUARTER_4      { \
  {SCORE_PLAY, LENGTH_QUARTER}, {SCORE_PLAY, LENGTH_QUARTER}, \
  {SCORE_PLAY, LENGTH_QUARTER}, {SCORE_PLAY, LENGTH_QUARTER}, \
  SCORE_FILL, SCORE_FILL, SCORE_FILL, SCORE_FILL \
}

#define BAR_QUARTER_3_REST_1 { \
  {SCORE_PLAY, LENGTH_QUARTER}, {SCORE_PLAY, LENGTH_QUARTER}, \
  {SCORE_PLAY, LENGTH_QUARTER}, {SCORE_REST, LENGTH_QUARTER}, \
  SCORE_FILL, SCORE_FILL, SCORE_FILL, SCORE_FILL \
}

#define BAR_EIGHTH_8       { \
  {SCORE_PLAY, LENGTH_EIGHTH}, {SCORE_PLAY, LENGTH_EIGHTH}, \
  {SCORE_PLAY, LENGTH_EIGHTH}, {SCORE_PLAY, LENGTH_EIGHTH}, \
  {SCORE_PLAY, LENGTH_EIGHTH}, {SCORE_PLAY, LENGTH_EIGHTH}, \
  {SCORE_PLAY, LENGTH_EIGHTH}, {SCORE_PLAY, LENGTH_EIGHTH} \
}

#define BAR_WHOLE_REST     { \
  {SCORE_REST, LENGTH_WHOLE}, \
  SCORE_FILL, SCORE_FILL, SCORE_FILL, SCORE_FILL, SCORE_FILL, SCORE_FILL, SCORE_FILL \
}

#define BAR_NOTE_REST_NOTE_REST { \
  {SCORE_PLAY, LENGTH_QUARTER}, {SCORE_REST, LENGTH_QUARTER}, \
  {SCORE_PLAY, LENGTH_QUARTER}, {SCORE_REST, LENGTH_QUARTER}, \
  SCORE_FILL, SCORE_FILL, SCORE_FILL, SCORE_FILL \
}

const ServoScoreEvent EntertainmentRobot::SCORE_TABLE[SERVER_BAR_COUNT][MAX_EVENTS_PER_BAR] = {
  BAR_QUARTER_4,            // 0小節目
  BAR_QUARTER_3_REST_1,     // 1小節目
  BAR_QUARTER_4,            // 2小節目
  BAR_QUARTER_3_REST_1,     // 3小節目
  BAR_NOTE_REST_NOTE_REST,  // 4小節目
  BAR_NOTE_REST_NOTE_REST,  // 5小節目
  BAR_EIGHTH_8,             // 6小節目
  BAR_QUARTER_3_REST_1,     // 7小節目
  
  BAR_WHOLE_REST,           // 8小節目
  BAR_WHOLE_REST,           // 9小節目
  BAR_WHOLE_REST,           // 10小節目
  BAR_WHOLE_REST,           // 11小節目
  BAR_WHOLE_REST,           // 12小節目

  BAR_QUARTER_4,            // 13小節目
  BAR_QUARTER_3_REST_1,     // 14小節目
  BAR_QUARTER_4,            // 15小節目
  BAR_QUARTER_3_REST_1,     // 16小節目
  BAR_NOTE_REST_NOTE_REST,  // 17小節目
  BAR_NOTE_REST_NOTE_REST,  // 18小節目
  BAR_EIGHTH_8,             // 19小節目
  BAR_QUARTER_3_REST_1,     // 20小節目

  BAR_WHOLE_REST,           // 21小節目
  BAR_WHOLE_REST,           // 22小節目
  BAR_WHOLE_REST,           // 23小節目
  BAR_WHOLE_REST            // 24小節目
};

#undef SCORE_FILL
#undef BAR_QUARTER_4
#undef BAR_QUARTER_3_REST_1
#undef BAR_EIGHTH_8
#undef BAR_WHOLE_REST
#undef BAR_NOTE_REST_NOTE_REST


EntertainmentRobot::EntertainmentRobot() :
    _currentBPM(120), _currentBar(0), _nextEventTime(0),
    _scoreEventIndex(0), _usedUnitsInBar(0),
    _barActive(false), _armDirection(false), _motionMutedByWholeRest(false),
    _scoreRestActive(true), _pianoUseRightArm(true),
    _rightArmReturnActive(false), _leftArmReturnActive(false), _extraReturnActive(false),
    _rightArmReturnTime(0), _leftArmReturnTime(0), _extraReturnTime(0),
    _emotionActive(false), _lastEmotionMoveTime(0), _emotionDirection(false)
{
    calculateBeatInterval();
}

void EntertainmentRobot::begin(const char* ssid, const char* pass, uint16_t port) {
    _localPort = port;
    
    setupServos();
    calculateBeatInterval();
    
    Serial.println("Connecting to Wi-Fi...");
    while (WiFi.begin(ssid, pass) != WL_CONNECTED) {
        Serial.print(".");
        delay(1000);
    }
    Serial.println("\nWi-Fi Connected!");

    _udp.begin(_localPort);

    Serial.println("Entertainment client started");
    Serial.print("IP: ");
    Serial.println(WiFi.localIP());
}

void EntertainmentRobot::update() {
    receiveUdpData();
    updateNormalMotion();
    updateEmotionMotion();
    updateServoReturn();
}

void EntertainmentRobot::setupServos() {
#if INSTRUMENT == PIANO
    _rightArmServo.attach(RIGHT_ARM_PIN);
    _leftArmServo.attach(LEFT_ARM_PIN);
    _rightArmServo.write(ARM_REST_ANGLE);
    _leftArmServo.write(ARM_REST_ANGLE);

#elif INSTRUMENT == TROMBONE
    _rightArmServo.attach(RIGHT_ARM_PIN);
    _rightArmServo.write(ARM_REST_ANGLE);

#elif INSTRUMENT == VIOLIN
    _rightArmServo.attach(RIGHT_ARM_PIN);
    _extraServo.attach(EXTRA_PIN);
    _rightArmServo.write(ARM_REST_ANGLE);
    _extraServo.write(HEAD_CENTER_ANGLE);

#elif INSTRUMENT == CASTANET
    _rightArmServo.attach(RIGHT_ARM_PIN);
    _extraServo.attach(EXTRA_PIN);
    _rightArmServo.write(ARM_REST_ANGLE);
    _extraServo.write(MOUTH_CLOSE_ANGLE);
#endif
}

void EntertainmentRobot::receiveUdpData() {
    int packetSize = _udp.parsePacket();

    if (packetSize > 0) {
        int data = _udp.read();

        while (_udp.available()) {
            _udp.read();
        }

        if (data < 0) {
            return;
        }

        // 40以上（または25以上）をBPMとして判定するための境界
        if (data >= 40) {
            updateBPM((uint8_t)data);
        } else {
            startBarMotion((uint8_t)data);
        }
    }
}

void EntertainmentRobot::updateBPM(uint8_t newBPM) {
    if (newBPM == _currentBPM) {
        return;
    }

    _currentBPM = newBPM;
    calculateBeatInterval();

    Serial.print("BPM updated: ");
    Serial.println(_currentBPM);

    startEmotion();
}

void EntertainmentRobot::calculateBeatInterval() {
    _beatInterval = 60000UL / _currentBPM;
}

unsigned long EntertainmentRobot::noteLengthToMs(uint8_t lengthUnits) {
    if (lengthUnits == 0) return 0;
    return (_beatInterval * (unsigned long)lengthUnits) / QUARTER_LENGTH_UNITS;
}

unsigned long EntertainmentRobot::calculatePulseTime(unsigned long eventDuration) {
    if (eventDuration == 0) return 0;
    unsigned long pulseTime = (eventDuration * 7UL) / 10UL;

    if (pulseTime < 40UL) { pulseTime = 40UL; }
    if (eventDuration > 30UL && pulseTime > eventDuration - 20UL) {
        pulseTime = eventDuration - 20UL;
    }
    if (pulseTime > eventDuration) { pulseTime = eventDuration; }

    return pulseTime;
}

uint8_t EntertainmentRobot::toScoreBar(uint8_t serverBarNumber) {
    int scoreBar = (int)serverBarNumber + PART_BAR_OFFSET;
    scoreBar %= SERVER_BAR_COUNT;

    if (scoreBar < 0) {
        scoreBar += SERVER_BAR_COUNT;
    }

    return (uint8_t)scoreBar;
}

bool EntertainmentRobot::isWholeRestBar(uint8_t scoreBarNumber) {
    if (scoreBarNumber >= SERVER_BAR_COUNT) {
        return false;
    }
    return WHOLE_REST_BAR_TABLE[scoreBarNumber];
}

void EntertainmentRobot::startBarMotion(uint8_t barNumber) {
    uint8_t scoreBar = toScoreBar(barNumber);
    _currentBar = scoreBar;

    Serial.print("Server bar: ");
    Serial.print(barNumber);
    Serial.print(" -> Score bar: ");
    Serial.println(_currentBar);

    if (isWholeRestBar(_currentBar)) {
        _motionMutedByWholeRest = true;
        _scoreRestActive = true;
        stopAllMotionForWholeRest();
        return;
    }

    _motionMutedByWholeRest = false;
    _scoreRestActive = false;

    _nextEventTime = millis();
    _scoreEventIndex = 0;
    _usedUnitsInBar = 0;
    _barActive = true;
}

void EntertainmentRobot::updateNormalMotion() {
    if (_motionMutedByWholeRest || !_barActive) {
        return;
    }

    while (_barActive && timeReached(_nextEventTime)) {
        triggerNextScoreEvent();
    }
}

void EntertainmentRobot::triggerNextScoreEvent() {
    if (_scoreEventIndex >= MAX_EVENTS_PER_BAR || _usedUnitsInBar >= BAR_LENGTH_UNITS) {
        finishBarMotion();
        return;
    }

    ServoScoreEvent event = SCORE_TABLE[_currentBar][_scoreEventIndex];
    _scoreEventIndex++;

    if (event.type == SCORE_END || event.lengthUnits == 0) {
        finishBarMotion();
        return;
    }

    uint8_t lengthUnits = event.lengthUnits;
    if (_usedUnitsInBar + lengthUnits > BAR_LENGTH_UNITS) {
        lengthUnits = BAR_LENGTH_UNITS - _usedUnitsInBar;
    }

    unsigned long eventDuration = noteLengthToMs(lengthUnits);

    if (event.type == SCORE_REST) {
        _scoreRestActive = true;
        stopMainMotionForRest();
    } else if (event.type == SCORE_PLAY) {
        _scoreRestActive = false;
        triggerScoreMotion(eventDuration);
    }

    _usedUnitsInBar += lengthUnits;
    _nextEventTime += eventDuration;
}

void EntertainmentRobot::triggerScoreMotion(unsigned long eventDuration) {
    unsigned long pulseTime = calculatePulseTime(eventDuration);

#if INSTRUMENT == PIANO
    if (_pianoUseRightArm) {
        triggerRightArmPulse(PIANO_RIGHT_PLAY_ANGLE, pulseTime);
    } else {
        triggerLeftArmPulse(PIANO_LEFT_PLAY_ANGLE, pulseTime);
    }
    _pianoUseRightArm = !_pianoUseRightArm;

#elif INSTRUMENT == TROMBONE
    _armDirection = !_armDirection;
    if (_armDirection) {
        _rightArmServo.write(TROMBONE_FORWARD_ANGLE);
    } else {
        _rightArmServo.write(TROMBONE_BACK_ANGLE);
    }

#elif INSTRUMENT == VIOLIN
    _armDirection = !_armDirection;
    if (_armDirection) {
        _rightArmServo.write(VIOLIN_BOW_LEFT_ANGLE);
    } else {
        _rightArmServo.write(VIOLIN_BOW_RIGHT_ANGLE);
    }

#elif INSTRUMENT == CASTANET
    triggerRightArmPulse(CASTANET_PLAY_ANGLE, pulseTime);
#endif
}

void EntertainmentRobot::finishBarMotion() {
    _barActive = false;
    _scoreRestActive = true;
    stopMainMotionForRest();
}

void EntertainmentRobot::startEmotion() {
    if (_motionMutedByWholeRest || _scoreRestActive) {
        return;
    }

#if INSTRUMENT == VIOLIN || INSTRUMENT == CASTANET
    _emotionActive = true;
    _emotionEndTime = millis() + _beatInterval * EMOTION_BEATS;
    _lastEmotionMoveTime = 0;
    _emotionDirection = false;
#endif
}

void EntertainmentRobot::updateEmotionMotion() {
    if (_motionMutedByWholeRest || _scoreRestActive) {
        resetEmotionServo();
        _emotionActive = false;
        return;
    }

    if (!_emotionActive) {
        return;
    }

    unsigned long now = millis();
    if (timeReached(_emotionEndTime)) {
        resetEmotionServo();
        _emotionActive = false;
        return;
    }

    unsigned long emotionInterval = _beatInterval / 2;
    if (emotionInterval < 100) {
        emotionInterval = 100;
    }

    if (now - _lastEmotionMoveTime >= emotionInterval) {
        _lastEmotionMoveTime = now;
        _emotionDirection = !_emotionDirection;

#if INSTRUMENT == VIOLIN
        if (_emotionDirection) {
            _extraServo.write(HEAD_LEFT_ANGLE);
        } else {
            _extraServo.write(HEAD_RIGHT_ANGLE);
        }
#elif INSTRUMENT == CASTANET
        if (_emotionDirection) {
            _extraServo.write(MOUTH_OPEN_ANGLE);
        } else {
            _extraServo.write(MOUTH_CLOSE_ANGLE);
        }
#endif
    }
}

void EntertainmentRobot::triggerRightArmPulse(int angle, unsigned long pulseTime) {
    _rightArmServo.write(angle);
    _rightArmReturnActive = true;
    _rightArmReturnTime = millis() + pulseTime;
}

void EntertainmentRobot::triggerLeftArmPulse(int angle, unsigned long pulseTime) {
    _leftArmServo.write(angle);
    _leftArmReturnActive = true;
    _leftArmReturnTime = millis() + pulseTime;
}

void EntertainmentRobot::triggerExtraPulse(int angle, unsigned long pulseTime) {
    _extraServo.write(angle);
    _extraReturnActive = true;
    _extraReturnTime = millis() + pulseTime;
}

void EntertainmentRobot::updateServoReturn() {
    unsigned long now = millis();

    if (_rightArmReturnActive && timeReached(_rightArmReturnTime)) {
        _rightArmServo.write(ARM_REST_ANGLE);
        _rightArmReturnActive = false;
    }

    if (_leftArmReturnActive && timeReached(_leftArmReturnTime)) {
        _leftArmServo.write(ARM_REST_ANGLE);
        _leftArmReturnActive = false;
    }

    if (_extraReturnActive && timeReached(_extraReturnTime)) {
        resetEmotionServo();
        _extraReturnActive = false;
    }
}

void EntertainmentRobot::stopAllMotionForWholeRest() {
    _barActive = false;
    _scoreEventIndex = 0;
    _usedUnitsInBar = 0;
    _emotionActive = false;
    cancelServoReturns();
    resetMainServos();
    resetEmotionServo();
}

void EntertainmentRobot::stopMainMotionForRest() {
    _rightArmReturnActive = false;
    _leftArmReturnActive = false;
    resetMainServos();
}

void EntertainmentRobot::cancelServoReturns() {
    _rightArmReturnActive = false;
    _leftArmReturnActive = false;
    _extraReturnActive = false;
}

void EntertainmentRobot::resetMainServos() {
#if INSTRUMENT == PIANO
    _rightArmServo.write(ARM_REST_ANGLE);
    _leftArmServo.write(ARM_REST_ANGLE);
#elif INSTRUMENT == TROMBONE
    _rightArmServo.write(ARM_REST_ANGLE);
#elif INSTRUMENT == VIOLIN
    _rightArmServo.write(ARM_REST_ANGLE);
#elif INSTRUMENT == CASTANET
    _rightArmServo.write(ARM_REST_ANGLE);
#endif
}

void EntertainmentRobot::resetEmotionServo() {
#if INSTRUMENT == VIOLIN
    _extraServo.write(HEAD_CENTER_ANGLE);
#elif INSTRUMENT == CASTANET
    _extraServo.write(MOUTH_CLOSE_ANGLE);
#endif
}

bool EntertainmentRobot::timeReached(unsigned long targetTime) {
    return (long)(millis() - targetTime) >= 0;
}
\end{lstlisting}

\subsubsection{カスタネットクライアント（castanets\_client.ino）}
カスタネット用のクライアントプログラムである．共通のクライアントプログラム（\texttt{client.ino}）を，音程を持たないカスタネット用に変更したものである．音階を並べたメロディの代わりに，固定周波数（1800\,Hz）の打撃と休符を組み合わせた打撃パターンを譜面データとして保持する．

\begin{lstlisting}
#include <WiFiS3.h>
#include <WiFiUdp.h>
#include "castanets_client_function.h"

// Offset変数とsetupScore関数を楽器に合わせて下さい．その他の，client_function.cpp，client_function.hはいじらなくても大丈夫です．

// --- ネットワーク設定 ---
char ssid[] = "hackathon003-WPA2"; //
char pass[] = "hackathon003"; //
uint16_t Port = 3000; // サーバーと同じポート番号

WiFiUDP Udp;

// --- クライアントサイドのグローバル変数 ---
char Offset = 0;         // 輪唱に必要なオフセット値（楽器ごとに設定，2番手だと-2）
uint8_t Data = 255;       // 受信データ
uint8_t CurrentBPM = 120; // 現在のBPM（初期値はサーバーに合わせて120とする）
uint8_t CurrentBar = 255; // 現在演奏中の小節番号
uint8_t NoteIndex = 255;  // 現在の小節の中で、何番目の音符を鳴らしているかのインデックス
bool Flag = false;        // 受信データが小節番号かBPMかを判定するフラグ
float ToneLength = 0.0;   // 基準音符（4分音符）の音の長さ（ミリ秒）
uint32_t StartTime = 0;   // 前回の音符を鳴らした時間（millis）
uint16_t Interval = 0;    // 実際に演奏する音の長さ（ミリ秒）

Pfm Score[40]; // 楽譜（40小節分の配列）

// --- 周波数定義 ---
constexpr int kFreqValue = 1800; // 音階なしのため固定の周波数
constexpr int kRest = 0;

void setupScore()
{
  // 第0小節〜第7小節（1周目）
  Score[0] = {{{kFreqValue, 0.2f}, {kRest, 0.8f}, {kRest, 1.0f}, {kFreqValue, 0.2f}, {kRest, 0.8f}, {kRest, 1.0f}}, 8};
  Score[1] = {{{kFreqValue, 0.2f}, {kRest, 0.8f}, {kRest, 1.0f}, {kFreqValue, 0.2f}, {kRest, 0.8f}, {kRest, 1.0f}}, 8};
  Score[2] = {{{kFreqValue, 0.2f}, {kRest, 0.8f}, {kRest, 1.0f}, {kFreqValue, 0.2f}, {kRest, 0.8f}, {kRest, 1.0f}}, 8};
  Score[3] = {{{kFreqValue, 0.2f}, {kRest, 0.8f}, {kRest, 1.0f}, {kFreqValue, 0.2f}, {kRest, 0.8f}, {kRest, 1.0f}}, 8};
  Score[4] = {{{kFreqValue, 0.2f}, {kRest, 0.8f}, {kRest, 1.0f}, {kFreqValue, 0.2f}, {kRest, 0.8f}, {kRest, 1.0f}}, 8};
  Score[5] = {{{kFreqValue, 0.2f}, {kRest, 0.8f}, {kRest, 1.0f}, {kFreqValue, 0.2f}, {kRest, 0.8f}, {kRest, 1.0f}}, 8};
  Score[6] = {{{kFreqValue, 0.2f}, {kRest, 0.8f}, {kRest, 1.0f}, {kFreqValue, 0.2f}, {kRest, 0.8f}, {kRest, 1.0f}}, 8};
  Score[7] = {{{kFreqValue, 0.2f}, {kRest, 0.8f}, {kRest, 1.0f}, {kFreqValue, 0.2f}, {kRest, 0.8f}, {kRest, 1.0f}}, 8};

  // 第8小節（休符）
  Score[8] = {{{kFreqValue, 0.2f}, {kRest, 0.8f}, {kRest, 1.0f}, {kFreqValue, 0.2f}, {kRest, 0.8f}, {kRest, 1.0f}}, 8};

  // 第9小節〜第16小節（2周目：1オクターブ上）
  Score[9] = {{{kFreqValue, 0.2f}, {kRest, 0.8f}, {kRest, 1.0f}, {kFreqValue, 0.2f}, {kRest, 0.8f}, {kRest, 1.0f}}, 8};
  Score[10] = {{{kFreqValue, 0.2f}, {kRest, 0.8f}, {kRest, 1.0f}, {kFreqValue, 0.2f}, {kRest, 0.8f}, {kRest, 1.0f}}, 8};
  Score[11] = {{{kFreqValue, 0.2f}, {kRest, 0.8f}, {kRest, 1.0f}, {kFreqValue, 0.2f}, {kRest, 0.8f}, {kRest, 1.0f}}, 8};
  Score[12] = {{{kFreqValue, 0.2f}, {kRest, 0.8f}, {kRest, 1.0f}, {kFreqValue, 0.2f}, {kRest, 0.8f}, {kRest, 1.0f}}, 8};
  Score[13] = {{{kFreqValue, 0.2f}, {kRest, 0.8f}, {kRest, 1.0f}, {kFreqValue, 0.2f}, {kRest, 0.8f}, {kRest, 1.0f}}, 8};
  Score[14] = {{{kFreqValue, 0.2f}, {kRest, 0.8f}, {kRest, 1.0f}, {kFreqValue, 0.2f}, {kRest, 0.8f}, {kRest, 1.0f}}, 8};
  Score[15] = {{{kFreqValue, 0.2f}, {kRest, 0.8f}, {kRest, 1.0f}, {kFreqValue, 0.2f}, {kRest, 0.8f}, {kRest, 1.0f}}, 8};
  Score[16] = {{{kFreqValue, 0.2f}, {kRest, 0.8f}, {kRest, 1.0f}, {kFreqValue, 0.2f}, {kRest, 0.8f}, {kRest, 1.0f}}, 8};

  // 第17小節〜第18小節（休符）
  Score[17] = {{{kFreqValue, 0.2f}, {kRest, 0.8f}, {kRest, 1.0f}, {kFreqValue, 0.2f}, {kRest, 0.8f}, {kRest, 1.0f}}, 8};
  Score[18] = {{{kFreqValue, 0.2f}, {kRest, 0.8f}, {kRest, 1.0f}, {kFreqValue, 0.2f}, {kRest, 0.8f}, {kRest, 1.0f}}, 8};

  // 第19小節〜第26小節（3周目：通常）
  Score[19] = {{{kFreqValue, 0.2f}, {kRest, 0.8f}, {kRest, 1.0f}, {kFreqValue, 0.2f}, {kRest, 0.8f}, {kRest, 1.0f}}, 8};
  Score[20] = {{{kFreqValue, 0.2f}, {kRest, 0.8f}, {kRest, 1.0f}, {kFreqValue, 0.2f}, {kRest, 0.8f}, {kRest, 1.0f}}, 8};
  Score[21] = {{{kFreqValue, 0.2f}, {kRest, 0.8f}, {kRest, 1.0f}, {kFreqValue, 0.2f}, {kRest, 0.8f}, {kRest, 1.0f}}, 8};
  Score[22] = {{{kFreqValue, 0.2f}, {kRest, 0.8f}, {kRest, 1.0f}, {kFreqValue, 0.2f}, {kRest, 0.8f}, {kRest, 1.0f}}, 8};
  Score[23] = {{{kFreqValue, 0.2f}, {kRest, 0.8f}, {kRest, 1.0f}, {kFreqValue, 0.2f}, {kRest, 0.8f}, {kRest, 1.0f}}, 8};
  Score[24] = {{{kFreqValue, 0.2f}, {kRest, 0.8f}, {kRest, 1.0f}, {kFreqValue, 0.2f}, {kRest, 0.8f}, {kRest, 1.0f}}, 8};
  Score[25] = {{{kFreqValue, 0.2f}, {kRest, 0.8f}, {kRest, 1.0f}, {kFreqValue, 0.2f}, {kRest, 0.8f}, {kRest, 1.0f}}, 8};
  Score[26] = {{{kFreqValue, 0.2f}, {kRest, 0.8f}, {kRest, 1.0f}, {kFreqValue, 0.2f}, {kRest, 0.8f}, {kRest, 1.0f}}, 8};

  // 第27小節（休符）
  Score[27] = {{{kFreqValue, 0.2f}, {kRest, 0.8f}, {kRest, 1.0f}, {kFreqValue, 0.2f}, {kRest, 0.8f}, {kRest, 1.0f}}, 8};

  // 第28小節〜第35小節（4周目：1オクターブ上）
  Score[28] = {{{kFreqValue, 0.2f}, {kRest, 0.8f}, {kRest, 1.0f}, {kFreqValue, 0.2f}, {kRest, 0.8f}, {kRest, 1.0f}}, 8};
  Score[29] = {{{kFreqValue, 0.2f}, {kRest, 0.8f}, {kRest, 1.0f}, {kFreqValue, 0.2f}, {kRest, 0.8f}, {kRest, 1.0f}}, 8};
  Score[30] = {{{kFreqValue, 0.2f}, {kRest, 0.8f}, {kRest, 1.0f}, {kFreqValue, 0.2f}, {kRest, 0.8f}, {kRest, 1.0f}}, 8};
  Score[31] = {{{kFreqValue, 0.2f}, {kRest, 0.8f}, {kRest, 1.0f}, {kFreqValue, 0.2f}, {kRest, 0.8f}, {kRest, 1.0f}}, 8};
  Score[32] = {{{kFreqValue, 0.2f}, {kRest, 0.8f}, {kRest, 1.0f}, {kFreqValue, 0.2f}, {kRest, 0.8f}, {kRest, 1.0f}}, 8};
  Score[33] = {{{kFreqValue, 0.2f}, {kRest, 0.8f}, {kRest, 1.0f}, {kFreqValue, 0.2f}, {kRest, 0.8f}, {kRest, 1.0f}}, 8};
  Score[34] = {{{kFreqValue, 0.2f}, {kRest, 0.8f}, {kRest, 1.0f}, {kFreqValue, 0.2f}, {kRest, 0.8f}, {kRest, 1.0f}}, 8};
  Score[35] = {{{kFreqValue, 0.2f}, {kRest, 0.8f}, {kRest, 1.0f}, {kFreqValue, 0.2f}, {kRest, 0.8f}, {kRest, 1.0f}}, 8};

  // 第36小節〜第39小節（終了・全休符）
  for (int i = 36; i < 40; i++)
  {
    Score[i] = {{{kRest, 4.0f}}, 1};
  }
}


void setup()
{
  Serial.begin(115200);

  Serial.println("Wi-Fiに接続中...");
  WiFi.begin(ssid, pass);
  while (WiFi.status() != WL_CONNECTED)
  {
    delay(500);
    Serial.print(".");
  }
  Serial.println("\nWi-Fi接続完了!");

  // IPアドレス取得待機
  while (WiFi.localIP() == IPAddress(0, 0, 0, 0))
  {
    Serial.println("ルーターからIPアドレスを取得中...");
    delay(500);
  }
  Serial.print("クライアントのIPアドレス: ");
  Serial.println(WiFi.localIP());

  Udp.begin(Port);

  setupScore();
  ToneLength = 60000.0 / CurrentBPM;
  Serial.println("クライアント：譜面データの読み込みが完了しました．");
}

void loop()
{
  // 1. 受信パケットの仕分けと小節番号のオフセット適用
  Flag = Parse_data(&Data, Offset, Udp);

  // 2. データがBPMだった場合の同期処理
  BPM_update(&CurrentBPM, Flag, Data, &ToneLength);

  // 3. 演奏位置制御とPCへのデータ送信（今回追加）
  Performance(Data, &CurrentBar, &NoteIndex, &StartTime, &Interval, ToneLength, Score);
}
\end{lstlisting}

\subsubsection{カスタネットクライアントの関数ヘッダ（castanets\_client\_function.h）}
カスタネット用の譜面データ構造体と関数のプロトタイプを定義したファイルである．カスタネットは音の強さを用いないため，共通版の\texttt{NoteData}構造体から強さ（\texttt{velocity}）を除いている．

\begin{lstlisting}
#ifndef CLIENT_FUNCTION_H
#define CLIENT_FUNCTION_H
#include <Arduino.h>
#include <WiFiUdp.h>

// --- 譜面データの構造体定義 ---
// 1つの音符の情報を表す構造体
struct NoteData {
  int pitch;      // 音の高さ（周波数 Hz）
  float length;   // 音の長さ（基準となる4分音符を 1.0 とした相対値）
};

// 1つの小節（Pfm）の情報を表す構造体
struct Pfm {
  NoteData notes[8]; // 1小節に最大8つの音符が入ると仮定
  int noteCount;     // この小節に実際に含まれる音符の数
};

bool Parse_data(uint8_t *Data, char Offset, WiFiUDP &Udp);// 受信データがBPM値ならtrue（フラグON），小節番号ならfalse（フラグOFF）を返す関数

void BPM_update(uint8_t *CurrentBPM, bool Flag, uint8_t Data, float *ToneLength);

void Performance(uint8_t Data, uint8_t *CurrentBar, uint8_t *NoteIndex, uint32_t *StartTime, uint16_t *Interval, float ToneLength, Pfm *Score);
#endif
\end{lstlisting}

\subsubsection{カスタネットクライアントの関数実装（castanets\_client\_function.cpp）}
カスタネット用のクライアントの処理を実装したファイルである．\texttt{Parse\_data}および\texttt{BPM\_update}の処理は共通版と同様である．\texttt{Performance}関数では，音の強さを送信せず，周波数と発音間隔のみをシリアル通信でProcessingへ送信する．

\begin{lstlisting}
#include "castanets_client_function.h"

bool Parse_data(uint8_t *Data, char Offset, WiFiUDP &Udp) {
  int packetSize = Udp.parsePacket();
  
  // 1. パケットを受信したかどうかの判定
  if (packetSize) {
    uint8_t received = Udp.read();
    Udp.flush();

    // === 【生存確認用デバッグ】受信した生の数値をそのままProcessingへ送る ===
    Serial.print("RAW_DATA:");
    Serial.println(received);
    // ===================================================================
    
    // 2. 受信データが40以上であればBPM値と判定（フラグONを返す）
    if (received >= 40) {
      *Data = received;
      return true; // フラグON
    }
    // 3. 40未満であれば小節番号と判定（オフセットを適用し，フラグOFFを返す）
    else {
      // サーバーの小節番号とオフセットを足す
      int myBar = received + Offset; 
      
      // マイナスになる場合（まだ自分の出番ではない待機期間）
      if (myBar < 0) {
        *Data = 255; // 待機状態を維持するための無効な値を入れる
      } else {
        *Data = myBar; // 自分の小節番号として適用
      }
      return false; // フラグOFF
    }
  }
  return false;
}

void BPM_update(uint8_t *CurrentBPM, bool Flag, uint8_t Data, float *ToneLength) {
  // フラグがOFF（小節番号）の場合は何もせず終了
  if (!Flag) {
    return;
  } 
  // フラグがON（BPM値）の場合の更新処理
  else {
    *CurrentBPM = Data;
    
    // 基準となる4分音符の長さをミリ秒単位で計算 (60000 / BPM)
    *ToneLength = 60000.0 / (*CurrentBPM);
    
    // 確認用のシリアル出力
    Serial.print("【BPM同期】新しいBPM: ");
    Serial.print(*CurrentBPM);
    Serial.print(" -> 4分音符の基準長さ(ms): ");
    Serial.println(*ToneLength);
  }
}

void Performance(uint8_t Data, uint8_t *CurrentBar, uint8_t *NoteIndex, uint32_t *StartTime, uint16_t *Interval, float ToneLength, Pfm *Score) {
  
  // 1. 新しい小節番号を受信した時のリセット処理
  // （Dataが40未満で、かつ現在演奏中の小節と異なる場合）
  if (Data < 40 && Data != *CurrentBar) {
    *CurrentBar = Data;    // 現在の小節番号を更新
    *NoteIndex = 0;        // 読み込む音符を最初(0番目)に戻す
    *StartTime = millis(); // 演奏開始タイマーをリセット
    *Interval = 0;         // 即座に最初の音が鳴るようにIntervalを0にする
  }

  // フリーズ防止
  if (*CurrentBar >= 40) {
    return; 
  }

  // 2. 音を鳴らすタイミング（時間が経過したか）の判定
  if (millis() - *StartTime >= *Interval) {
    
    // その小節にある音符の数(noteCount)に達するまで処理を行う
    if (*NoteIndex < Score[*CurrentBar].noteCount) {
      
      *StartTime = millis(); // 次の音の基準時間を現在時刻に更新
      
      // 譜面配列から「音の高さ(pitch)」と「音の長さの割合(length)」を取得
      int pitch = Score[*CurrentBar].notes[*NoteIndex].pitch;
      float length = Score[*CurrentBar].notes[*NoteIndex].length;
      
      // 次の音までの間隔（Interval）を計算（ミリ秒）
      *Interval = ToneLength * length;
      
      // PC（Processing）へシリアル通信で送信
      // 形式：「周波数,鳴らすミリ秒」 (例： 262,500)
      Serial.print(pitch);
      Serial.print(",");
      Serial.println(*Interval);
      
      *NoteIndex = *NoteIndex + 1; // 次の音符へ進む
    }
  }
}
\end{lstlisting}

\subsubsection{ピアノの音源（piano.pde）}
木村が担当したピアノの音源のソースコードである．\texttt{PianoSound}クラスが音源の中心であり，基本となる音，高域の輝きを与える音，および余韻を表す音の3種類のオシレータ群を重ね合わせることで，ピアノの音色を合成する．\texttt{serialEvent()}関数でクライアントから受信した音符の情報を解釈して発音する．

\begin{lstlisting}
import processing.serial.*;
import ddf.minim.*;
import ddf.minim.ugens.*;

// ==================================================
// シリアル通信
// ==================================================
Serial myPort;

// 前に確認したArduinoのポート番号
// [2] "/dev/cu.usbmodem34B7DA6537F82"
final int PORT_INDEX = 2;

// Arduino側が Serial.begin(115200); なので合わせる
final int BAUD_RATE = 115200;

// ==================================================
// Minim
// ==================================================
Minim minim;
AudioOutput out;

// ==================================================
// 現在の状態
// ==================================================
String currentNote = "REST";
float currentFreq = 0;
int currentDurationMs = 500;
int currentBpm = 120;
String lastReceived = "";

// ==================================================
// テスト再生用
// ==================================================
boolean testPlaying = false;
int testIndex = 0;
int nextNoteTime = 0;

float[] testFreqs = {
  262, 294, 330, 349,
  330, 294, 262, 0,
  330, 349, 392, 440,
  392, 349, 330, 0,
  262, 262, 262, 262,
  262, 294, 330, 349,
  330, 294, 262, 0
};


// ==================================================
// setup
// ==================================================
void setup() {
  size(600, 300);

  minim = new Minim(this);
  out = minim.getLineOut(Minim.MONO, 2048);

  println("=== Serial Port List ===");
  printArray(Serial.list());

  if (Serial.list().length > PORT_INDEX) {
    String portName = Serial.list()[PORT_INDEX];
    println("接続ポート: " + portName);

    myPort = new Serial(this, portName, BAUD_RATE);
    myPort.bufferUntil('\n');
    myPort.clear();
  } else {
    println("注意: PORT_INDEX が範囲外です．");
    println("Serial.list() の番号を確認してください．");
    println("pキーによるProcessing単体のテスト再生はできます．");
  }

  textSize(24);

  println("Processing起動完了");
  println("受信形式：周波数,長さms");
  println("例：262,500");
  println("休符：0,500");
  println("pキー：Processing単体でテスト再生");
}


// ==================================================
// draw
// ==================================================
void draw() {
  background(255);

  fill(0);
  textSize(24);
  text("Piano Part", 230, 70);

  textSize(18);
  text("Press P to test", 220, 115);
  text("Note: " + currentNote, 220, 155);
  text("Freq: " + currentFreq + " Hz", 220, 185);
  text("Duration: " + currentDurationMs + " ms", 180, 215);
  text("BPM: " + currentBpm, 245, 245);

  textSize(12);
  text("Last: " + lastReceived, 80, 280);

  // Pキーによるテスト再生
  if (testPlaying && millis() >= nextNoteTime) {
    playTestNote();

    float beatLength = 60.0 / currentBpm;
    nextNoteTime = millis() + int(beatLength * 1000);

    testIndex++;
  }

  if (testPlaying && testIndex >= testFreqs.length) {
    testPlaying = false;
    currentNote = "REST";
    currentFreq = 0;
  }
}


// ==================================================
// keyPressed
// ==================================================
void keyPressed() {
  if (key == 'p' || key == 'P') {
    testPlaying = true;
    testIndex = 0;
    nextNoteTime = millis();
    println("テスト再生開始");
  }
}


// ==================================================
// テスト再生
// ==================================================
void playTestNote() {
  if (testIndex < testFreqs.length) {
    float freq = testFreqs[testIndex];

    currentFreq = freq;
    currentNote = freqToName(freq);

    if (freq > 0) {
      float beatLength = 60.0 / currentBpm;

      // 透明感ある鍵盤音版
      out.playNote(0, beatLength, new PianoSound(freq, 0.50, beatLength));
    }
  }
}


// ==================================================
// Arduinoから受信したときに呼ばれる
//
// 今のArduino側の基本形式：
// 周波数,長さms
//
// 例：
// 262,500
// 294,500
// 0,500
//
// 0は休符
// ==================================================
void serialEvent(Serial p) {
  String data = p.readStringUntil('\n');

  if (data == null) {
    return;
  }

  data = trim(data);

  if (data.length() == 0) {
    return;
  }

  lastReceived = data;
  println("受信: " + data);

  // Arduino側のデバッグ出力は演奏データではないので無視
  if (data.startsWith("RAW_DATA:")) {
    println("デバッグデータなので無視: " + data);
    return;
  }

  // Wi-Fi接続メッセージなど，カンマがない行は無視
  if (data.indexOf(",") == -1) {
    println("演奏データではないため無視: " + data);
    return;
  }

  String[] values = split(data, ',');

  // ==================================================
  // 今のArduino側の形式：周波数,長さms
  // 例：262,500
  // ==================================================
  if (values.length == 2) {
    float freq = 0;
    int durationMs = 500;

    try {
      freq = float(trim(values[0]));
    }
    catch (Exception e) {
      println("周波数の読み取りに失敗: " + values[0]);
      return;
    }

    try {
      durationMs = int(trim(values[1]));
    }
    catch (Exception e) {
      println("長さmsの読み取りに失敗: " + values[1]);
      durationMs = 500;
    }

    currentFreq = freq;
    currentDurationMs = durationMs;
    currentNote = freqToName(freq);

    playFrequency(freq, durationMs);
    return;
  }

  // ==================================================
  // 将来用：周波数,長さms,BPM
  // 例：262,500,120
  // ==================================================
  if (values.length >= 3) {
    float freq = 0;
    int durationMs = 500;
    int bpm = currentBpm;

    try {
      freq = float(trim(values[0]));
    }
    catch (Exception e) {
      println("周波数の読み取りに失敗: " + values[0]);
      return;
    }

    try {
      durationMs = int(trim(values[1]));
    }
    catch (Exception e) {
      println("長さmsの読み取りに失敗: " + values[1]);
      durationMs = 500;
    }

    try {
      bpm = int(trim(values[2]));
    }
    catch (Exception e) {
      println("BPMの読み取りに失敗: " + values[2]);
      bpm = currentBpm;
    }

    currentFreq = freq;
    currentDurationMs = durationMs;
    currentBpm = bpm;
    currentNote = freqToName(freq);

    playFrequency(freq, durationMs);
    return;
  }

  println("形式が不明なため無視: " + data);
}


// ==================================================
// 周波数を鳴らす
// ==================================================
void playFrequency(float freq, int durationMs) {
  if (freq <= 0) {
    println("休符");
    currentNote = "REST";
    return;
  }

  float durationSec = durationMs / 1000.0;

  println("再生: " + freq + " Hz / " + durationMs + " ms");

  // 透明感ある鍵盤音版
  out.playNote(0, durationSec, new PianoSound(freq, 0.50, durationSec));
}




// ==================================================
// 周波数を音名に変換
// ==================================================
String freqToName(float freq) {
  if (freq <= 0) {
    return "REST";
  }

  if (abs(freq - 262) < 3) return "C4";
  if (abs(freq - 277) < 3) return "C#4";
  if (abs(freq - 294) < 3) return "D4";
  if (abs(freq - 311) < 3) return "D#4";
  if (abs(freq - 330) < 3) return "E4";
  if (abs(freq - 349) < 3) return "F4";
  if (abs(freq - 370) < 3) return "F#4";
  if (abs(freq - 392) < 3) return "G4";
  if (abs(freq - 415) < 3) return "G#4";
  if (abs(freq - 440) < 3) return "A4";
  if (abs(freq - 466) < 3) return "A#4";
  if (abs(freq - 494) < 3) return "B4";
  if (abs(freq - 523) < 3) return "C5";
  if (abs(freq - 587) < 3) return "D5";
  if (abs(freq - 659) < 3) return "E5";
  if (abs(freq - 698) < 3) return "F5";
  if (abs(freq - 784) < 3) return "G5";
  if (abs(freq - 880) < 3) return "A5";

  return nf(freq, 0, 1) + "Hz";
}

// ==================================================
// Piano sound ピアノ感強化版 V2
//
// 前回より大きく変更した点：
// ・整数倍音ではなく，ピアノ弦っぽい少しズレた倍音にする
// ・打鍵直後だけ「コツン」と鳴る高い成分を足す
// ・高い倍音ほど急速に消えるようにする
// ・木の箱っぽい共鳴成分を少しだけ足す
//
// 注意：Processing の合成音だけなので，本物の録音サンプルほどにはならない．
// ただし，前回より「ピアノらしい減衰」と「打鍵感」は強くなる．
// ==================================================
class PianoSound implements Instrument {
  Oscil[] toneWaves;
  ADSR[] toneEnvs;

  Oscil[] shineWaves;
  ADSR[] shineEnvs;

  Oscil[] tailWaves;
  ADSR[] tailEnvs;

  float freq;
  float amp;
  float durationSec;

  final int NUM_TONES = 10;
  final int NUM_SHINE = 5;
  final int NUM_TAIL = 3;

  PianoSound(float freq, float amp, float durationSec) {
    this.freq = freq;
    this.amp = amp;
    this.durationSec = durationSec;

    toneWaves = new Oscil[NUM_TONES];
    toneEnvs = new ADSR[NUM_TONES];

    shineWaves = new Oscil[NUM_SHINE];
    shineEnvs = new ADSR[NUM_SHINE];

    tailWaves = new Oscil[NUM_TAIL];
    tailEnvs = new ADSR[NUM_TAIL];

    // --------------------------------------------------
    // 透明感ある鍵盤音版
    // --------------------------------------------------
    // 前回の「ハンマー感」「箱鳴り」を弱める．
    // 機械音っぽさを減らすため，音の立ち上がりを少し丸くし，
    // 高い倍音は薄く短く入れる．

    // 全体音量．大きいとすぐ音割れして機械音になるので控えめ．
    float noteAmp = amp * 0.44;
    if (freq >= 500) noteAmp *= 0.88;
    if (freq >= 700) noteAmp *= 0.82;

    // 高音ほど高倍音を減らす．キンキン感を抑えるため．
    float highScale = map(freq, 262, 880, 0.88, 0.28);
    highScale = constrain(highScale, 0.28, 0.88);

    // 低音は少しだけ太くする．ただし濁らない程度．
    float lowScale = map(freq, 262, 880, 1.08, 0.82);
    lowScale = constrain(lowScale, 0.82, 1.08);

    // 3本弦っぽい薄い揺れ + 少しだけズレた倍音．
    // 完全な整数倍音だと電子音っぽくなるので，ほんの少しだけズラす．
    float[] toneRatio = {
      0.9992,
      1.0000,
      1.0008,
      2.004,
      3.010,
      4.018,
      5.030,
      6.045,
      8.070,
      10.100
    };

    // メイン成分．
    // 高い倍音を入れすぎないことで透明感を出す．
    float[] toneAmp = {
      noteAmp * 0.20 * lowScale,
      noteAmp * 0.34 * lowScale,
      noteAmp * 0.18 * lowScale,
      noteAmp * 0.100 * highScale,
      noteAmp * 0.055 * highScale,
      noteAmp * 0.030 * highScale,
      noteAmp * 0.018 * highScale,
      noteAmp * 0.011 * highScale,
      noteAmp * 0.006 * highScale,
      noteAmp * 0.0035 * highScale
    };

    // 打鍵音を丸くするため，Attackは前回より少し遅め．
    float attackTime = map(freq, 262, 880, 0.008, 0.004);
    attackTime = constrain(attackTime, 0.004, 0.008);

    // 透明感を残すため，低い成分は少し長め，高い成分は短め．
    float baseDecay = constrain(durationSec * 1.05, 0.22, 0.95);

    float[] toneDecay = {
      baseDecay * 1.12,
      baseDecay * 1.22,
      baseDecay * 1.08,
      baseDecay * 0.52,
      baseDecay * 0.36,
      baseDecay * 0.25,
      baseDecay * 0.18,
      baseDecay * 0.13,
      baseDecay * 0.09,
      baseDecay * 0.065
    };

    float[] toneSustain = {
      0.0030,
      0.0035,
      0.0028,
      0.00070,
      0.00045,
      0.00028,
      0.00018,
      0.00010,
      0.00006,
      0.00000
    };

    // 余韻．短すぎると電子音のブツ切れ感が出るので少し長め．
    float releaseTime = map(freq, 262, 880, 0.34, 0.14);
    releaseTime = constrain(releaseTime, 0.14, 0.34);

    for (int i = 0; i < NUM_TONES; i++) {
      // 透明感重視なので基本はSINEだけ．TRIANGLEは機械音感が出やすいので使わない．
      toneWaves[i] = new Oscil(freq * toneRatio[i], 1.0, Waves.SINE);

      toneEnvs[i] = new ADSR(
        toneAmp[i],
        attackTime,
        toneDecay[i],
        toneSustain[i],
        releaseTime
      );

      toneWaves[i].patch(toneEnvs[i]);
    }

    // --------------------------------------------------
    // きらっとした透明感成分
    // --------------------------------------------------
    // 打鍵音ではなく，ガラスっぽい薄い響き．
    // 強すぎると鉄琴になるのでかなり小さくする．
    float shineScale = map(freq, 262, 880, 0.70, 0.24);
    shineScale = constrain(shineScale, 0.24, 0.70);

    float[] shineRatio = {
      2.414,
      3.018,
      5.027,
      7.080,
      11.120
    };

    float[] shineAmp = {
      noteAmp * 0.013 * shineScale,
      noteAmp * 0.010 * shineScale,
      noteAmp * 0.0065 * shineScale,
      noteAmp * 0.0036 * shineScale,
      noteAmp * 0.0019 * shineScale
    };

    float[] shineDecay = {
      0.115,
      0.090,
      0.065,
      0.045,
      0.032
    };

    for (int i = 0; i < NUM_SHINE; i++) {
      shineWaves[i] = new Oscil(freq * shineRatio[i], 1.0, Waves.SINE);
      shineEnvs[i] = new ADSR(
        shineAmp[i],
        0.006,
        shineDecay[i],
        0.0,
        0.055
      );
      shineWaves[i].patch(shineEnvs[i]);
    }

    // --------------------------------------------------
    // 薄い余韻成分
    // --------------------------------------------------
    // リバーブの代わりに，かなり小さい音量で長く残る成分を足す．
    // これで音の終わりが少し空気っぽくなる．
    float tailScale = map(freq, 262, 880, 0.75, 0.35);
    tailScale = constrain(tailScale, 0.35, 0.75);

    float[] tailRatio = {
      1.0000,
      2.0020,
      3.0060
    };

    float[] tailAmp = {
      noteAmp * 0.030 * tailScale,
      noteAmp * 0.012 * tailScale,
      noteAmp * 0.006 * tailScale
    };

    for (int i = 0; i < NUM_TAIL; i++) {
      tailWaves[i] = new Oscil(freq * tailRatio[i], 1.0, Waves.SINE);
      tailEnvs[i] = new ADSR(
        tailAmp[i],
        0.020,
        constrain(durationSec * 1.20, 0.30, 1.10),
        0.0,
        releaseTime * 1.25
      );
      tailWaves[i].patch(tailEnvs[i]);
    }
  }

  void noteOn(float dur) {
    for (int i = 0; i < NUM_TONES; i++) {
      toneEnvs[i].patch(out);
      toneEnvs[i].noteOn();
    }

    for (int i = 0; i < NUM_SHINE; i++) {
      shineEnvs[i].patch(out);
      shineEnvs[i].noteOn();
    }

    for (int i = 0; i < NUM_TAIL; i++) {
      tailEnvs[i].patch(out);
      tailEnvs[i].noteOn();
    }
  }

  void noteOff() {
    for (int i = 0; i < NUM_TONES; i++) {
      toneEnvs[i].noteOff();
      toneEnvs[i].unpatchAfterRelease(out);
    }

    for (int i = 0; i < NUM_SHINE; i++) {
      shineEnvs[i].noteOff();
      shineEnvs[i].unpatchAfterRelease(out);
    }

    for (int i = 0; i < NUM_TAIL; i++) {
      tailEnvs[i].noteOff();
      tailEnvs[i].unpatchAfterRelease(out);
    }
  }
}
\end{lstlisting}

\subsubsection{ヴァイオリンの音源（violin.pde）}
仲里が担当したヴァイオリンの音源のソースコードである．\texttt{ViolinInstrument}クラスが音源の中心であり，胴の響きを表す波形，弦の成分を表す2つののこぎり波，およびビブラートを与えるオシレータを組み合わせることで，ヴァイオリンの音色を合成する．\texttt{serialEvent()}関数でクライアントから受信した音符の情報を解釈して発音する．

\begin{lstlisting}
// =====================================================================
// 「かえるのうた」 ヴァイオリン・パート（Processing側）― 最小表示版 v4
// ---------------------------------------------------------------------
// クライアントArduino（violin_client_function.cpp の Performance()）から
//   "周波数,鳴らすミリ秒,弾き方\n"   例） 392.00,500,0
// を受信し，ヴァイオリン音色で演奏します。
//
//   弾き方の値（Arduino から届く数値）：
//     0 = 下げ弓（しっかりした弓のストローク）
//     1 = 上げ弓（軽い弓のストローク）   ← 0と1の交互で「弓の上下」を表現
//     2 = デタシェ（テスト用）
//   ※「クヮ」も「ケケケ」も，譜面側で 0/1 を交互に並べてあるので，
//     下げ弓・上げ弓が交互に鳴り，弓を上下に返す質感になります。
//   ※音符の長さに応じて立ち上がりの鋭さを自動調整（速い刻みは鋭く，
//     ゆっくりした旋律は表情豊かに）。
//
// 【必要なライブラリ】 Minim
// 【セットアップ】 portName を自分の環境のポート名に書き換える。
// 【Arduino無しのテスト】
//   1〜6 = ド〜ラ（押すたびに下げ弓→上げ弓を交互），0 = 休符，
//   t = 「ケケケケ」走句（下げ/上げ交互）を自動で鳴らす。
// =====================================================================

import processing.serial.*;
import ddf.minim.*;
import ddf.minim.ugens.*;

Minim minim;
AudioOutput out;
Serial port;

Wavetable customViolinWave;

// ★自分の環境のポート名に書き換えてください
String portName = "/dev/cu.usbmodem34B7DA6377FC2";

// ---- 表示用の状態 ----
String  currentNoteName    = "ーー";
boolean currentIsRest      = true;
boolean hasReceivedAnyData = false;

// ---- 受信ログ（発音数の検証用CSV） ----
//   起動時に violin_recv_log.csv を作成し、受信した音符を1行ずつ記録する。
//   k キーで保存して閉じる（閉じなくても1行ごとに書き込まれる）。
PrintWriter recvLog;
int recvCount = 0;

// =====================================================================
// ★音色つまみ（演奏を聞きながらキーで調整できる）
//   金管/ポォー → 弦らしく したいときの目安：
//     ・暗くこもる(ポォー)  → gBright を上げる(q)
//     ・弦の擦れ感が欲しい   → gBowNoise を上げる(w)
//     ・金管っぽい鳴き       → gReso を下げる(d)
//     ・丸い/ホルンっぽい    → gSaw を上げ gBody を下げる(r)
//   キー： q/a=明るさ  w/s=擦れ音  e/d=共振  r/f=ノコギリ⇔胴
//   現在値はコンソールと画面左上に表示。良い値が決まったら教えてください。
// =====================================================================
float gBright   = 4000;   // フィルタ明るさ(Hz)。低いと倍音が消えて純音(ポォー)化する
float gReso     = 0.0;    // フィルタ共振。高い=金管っぽい鳴き／0=素直
float gBowNoise = 0.22;   // 弓の擦れ音の量。高い=弦らしいザラつき
float gSaw      = 0.60;   // ノコギリ波(弦成分=倍音が豊か)の量
float gBody     = 0.20;   // 胴成分(丸み・純音寄り)の量

void setup() {
  size(800, 400);

  minim = new Minim(this);
  out = minim.getLineOut();

  PFont jpFont = findJapaneseFont(28);
  if (jpFont != null) textFont(jpFont);

  // ヴァイオリンらしい倍音波形（violin.pde と同じ作り方）
  float[] harmonics = {0.8729, 1.0000, 0.1919, 0.1815, 0.1828, 0.0367, 0.0338, 0.0533, 0.0968, 0.0600};
  float[] waveData = new float[1024];
  float maxVal = 0;
  for (int i = 0; i < 1024; i++) {
    float t = (float)i / 1024.0 * TWO_PI;
    float val = 0;
    for (int h = 0; h < harmonics.length; h++) val += harmonics[h] * sin((h + 1) * t);
    waveData[i] = val;
    if (abs(val) > maxVal) maxVal = abs(val);
  }
  for (int i = 0; i < 1024; i++) waveData[i] /= maxVal;
  customViolinWave = new Wavetable(waveData);

  // 受信ログCSVを作成（ヘッダ行を書く）
  recvLog = createWriter("violin_recv_log.csv");
  recvLog.println("番号,周波数Hz,長さms,弾き方,音名,種別");
  recvLog.flush();

  println("利用可能なシリアルポート一覧:");
  printArray(Serial.list());
  openSerialPort();
  println("「かえるのうた」ヴァイオリン・パート 準備完了。");
}

void openSerialPort() {
  try {
    port = new Serial(this, portName, 115200);
  } catch (RuntimeException e) {
    println("ポート [" + portName + "] を開けませんでした: " + e.getMessage());
    String[] avail = Serial.list();
    if (avail.length > 0) {
      println("代わりに先頭のポート [" + avail[0] + "] を使用します。");
      port = new Serial(this, avail[0], 115200);
    } else {
      println("利用可能なシリアルポートが見つかりません。");
      return;
    }
  }
  port.bufferUntil('\n');
}

PFont findJapaneseFont(int fontSize) {
  String[] candidates = {
    "Hiragino Sans", "Hiragino Kaku Gothic ProN", "Yu Gothic", "YuGothic",
    "Meiryo", "MS Gothic", "MS PGothic", "Noto Sans CJK JP", "Noto Sans JP",
    "Source Han Sans JP", "IPAGothic", "TakaoGothic"
  };
  String[] available = PFont.list();
  for (String cand : candidates) {
    for (String avail : available) {
      if (avail.toLowerCase().contains(cand.toLowerCase())) {
        return createFont(avail, fontSize, true);
      }
    }
  }
  return null;
}

// ==========================================
// 描画：波形（上半分）＋ 音階（下半分）だけの最小構成
// ==========================================
void draw() {
  background(20, 30, 40);

  updateTestRun(); // テスト走句スケジューラ

  // --- 1. 波形（オシロスコープ）：中心Y=100 ---
  stroke(150, 255, 150);
  strokeWeight(2);
  noFill();
  beginShape();
  for (int i = 0; i < out.bufferSize(); i++) {
    float x = map(i, 0, out.bufferSize() - 1, 0, width);
    vertex(x, 100 + out.left.get(i) * 100);
  }
  endShape();
  strokeWeight(1);

  // --- 2. 音階：画面下半分に大きく表示 ---
  textAlign(CENTER, CENTER);
  fill(currentIsRest ? color(120, 140, 130) : color(255, 244, 214));
  textSize(120);
  String shown = hasReceivedAnyData ? currentNoteName : "ーー";
  text(shown, width / 2, 290);

  // --- 3. 音色つまみの現在値（左上）---
  textAlign(LEFT, TOP);
  fill(180, 200, 220);
  textSize(13);
  text("q/a 明るさ:" + int(gBright)
     + "  w/s 擦れ:" + nf(gBowNoise,1,2)
     + "  e/d 共振:" + nf(gReso,1,2)
     + "  r/f ノコギリ:" + nf(gSaw,1,2) + "/胴:" + nf(gBody,1,2),
       10, 8);
}

// ==========================================
// シリアル通信：3つのデータを受け取る
// ==========================================
void serialEvent(Serial p) {
  String inString = p.readStringUntil('\n');
  if (inString == null) return;
  inString = trim(inString);
  String[] data = split(inString, ',');
  if (data.length != 3) return; // RAW_DATA: 等のデバッグ行は無視

  float freq = float(data[0]);
  float durationMs = float(data[1]);
  int articulation = int(data[2]);

  // 受信した音符をCSVに記録（発音数の検証用）
  recvCount++;
  String shubetsu = (freq > 0) ? "発音" : "休符";
  recvLog.println(recvCount + "," + int(freq) + "," + int(durationMs) + ","
                  + articulation + "," + noteNameFor(freq) + "," + shubetsu);
  recvLog.flush();

  triggerNote(freq, durationMs, articulation);
}

void triggerNote(float freq, float durationMs, int articulation) {
  currentNoteName = noteNameFor(freq);
  currentIsRest = (freq <= 0);
  hasReceivedAnyData = true;

  if (freq > 0) {
    float durationSec = durationMs / 1000.0;
    float playDuration = durationSec * 0.85;
    if (playDuration < 0.05) playDuration = 0.05;
    // 音符の長さを合成側に渡し，立ち上がりの鋭さを自動調整
    out.playNote(0, playDuration, new ViolinInstrument(freq, 0.8, articulation, playDuration));
  }
}

String noteNameFor(float freq) {
  if (freq <= 0) return "休符";
  float[]  freqs = {262, 294, 330, 349, 392, 440, 523, 587, 659, 698, 784, 880};
  String[] names = {"ド", "レ", "ミ", "ファ", "ソ", "ラ", "高ド", "高レ", "高ミ", "高ファ", "高ソ", "高ラ"};
  int best = 0;
  float bestDiff = 999999.0;
  for (int i = 0; i < freqs.length; i++) {
    float diff = abs(freq - freqs[i]);
    if (diff < bestDiff) { bestDiff = diff; best = i; }
  }
  return names[best];
}

// ---- テスト用 ----
int   testBow = 0; // 0=下げ弓, 1=上げ弓 を交互に
float[] testRunFreqs;
int     testRunIndex = -1;
int     testRunNextMs = 0;
int     testRunStepMs = 220;

void keyPressed() {
  float[] tf = {262, 294, 330, 349, 392, 440};
  if (key >= '1' && key <= '6') {
    triggerNote(tf[key - '1'], 500, testBow);
    testBow = 1 - testBow;
  } else if (key == '0') {
    triggerNote(0, 500, 0);
  } else if (key == 't' || key == 'T') {
    testRunFreqs = new float[]{262, 262, 294, 294, 330, 330, 349, 349};
    testRunIndex = 0;
    testRunNextMs = millis();
  }
  // ---- 音色つまみ（聞きながら調整）----
  else if (key == 'q') { gBright += 200;                 printTone(); }
  else if (key == 'a') { gBright = max(500, gBright-200); printTone(); }
  else if (key == 'w') { gBowNoise += 0.03;              printTone(); }
  else if (key == 's') { gBowNoise = max(0, gBowNoise-0.03); printTone(); }
  else if (key == 'e') { gReso = min(0.6, gReso+0.03);   printTone(); }
  else if (key == 'd') { gReso = max(0, gReso-0.03);     printTone(); }
  else if (key == 'r') { gSaw += 0.05; gBody = max(0, gBody-0.05); printTone(); } // 弦寄り
  else if (key == 'f') { gSaw = max(0, gSaw-0.05); gBody += 0.05;  printTone(); } // 胴/丸み寄り
  // ---- 受信ログの保存 ----
  else if (key == 'k' || key == 'K') {
    recvLog.flush();
    println("受信ログ保存: violin_recv_log.csv（" + recvCount + " 行）");
  }
}

void printTone() {
  println("[音色] 明るさ gBright=" + int(gBright)
        + "  共振 gReso=" + nf(gReso,1,2)
        + "  擦れ gBowNoise=" + nf(gBowNoise,1,2)
        + "  ノコギリ gSaw=" + nf(gSaw,1,2)
        + "  胴 gBody=" + nf(gBody,1,2));
}

void updateTestRun() {
  if (testRunIndex < 0) return;
  if (millis() >= testRunNextMs) {
    if (testRunIndex < testRunFreqs.length) {
      // 下げ弓(0)→上げ弓(1)を交互に（実際の譜面と同じ）
      triggerNote(testRunFreqs[testRunIndex], testRunStepMs, testRunIndex % 2);
      testRunIndex++;
      testRunNextMs += testRunStepMs;
    } else {
      testRunIndex = -1;
    }
  }
}

void stop() {
  if (recvLog != null) { recvLog.flush(); recvLog.close(); }
  out.close();
  minim.stop();
  super.stop();
}

// ==========================================
// ヴァイオリンクラス（v4）
//   ・ノコギリ波の太さを復活（弦の胴鳴り重視）
//   ・低い音（クヮのC4等）はフィルターを暖かい帯域まで下げて金管化を回避
//   ・弓の擦れ音（PINKノイズ）を全弓使いで維持して弦の質感を確保
//   ・0=下げ弓 / 1=上げ弓 を弾き分け，交互で「弓の上下」を表現
//   ・音符の長さに応じて立ち上がり/余韻を自動調整
// ==========================================
class ViolinInstrument implements Instrument {
  Oscil waveBody, waveString1, waveString2, vibrato;
  Noise bowAttack, bowFriction;
  ADSR attackEnv, frictionEnv, ampEnv;
  Summer mix;
  MoogFilter filter;
  Line pitchSlide;

  ViolinInstrument(float frequency, float amplitude, int art, float noteLenSec) {

    // 短い音ほど立ち上がり・余韻を詰める係数（速い刻みは鋭く，旋律は豊かに）
    float lenF = constrain(noteLenSec / 0.45, 0.40, 1.25);

    // 弦成分(ノコギリ gSaw) ＋ 胴成分(丸み gBody)。つまみで比率を変えられる
    waveBody    = new Oscil(frequency,         amplitude * gBody,      customViolinWave);
    waveString1 = new Oscil(frequency * 1.002, amplitude * gSaw * 0.5, Waves.SAW);
    waveString2 = new Oscil(frequency * 0.998, amplitude * gSaw * 0.5, Waves.SAW);

    // ピッチのしゃくり（下げ弓は強め，上げ弓は軽め）
    float slideTime, slideStart;
    if (art == 1) {            // 上げ弓
      slideTime  = 0.05 * lenF; slideStart = frequency - 7;
    } else if (art == 2) {     // デタシェ
      slideTime  = 0.03 * lenF; slideStart = frequency - 5;
    } else {                   // 下げ弓
      slideTime  = 0.09 * lenF; slideStart = frequency - 11;
    }
    pitchSlide = new Line(slideTime, slideStart, frequency);

    vibrato = new Oscil(5.5, frequency * 0.012, Waves.SINE);
    pitchSlide.patch(vibrato.offset);
    vibrato.patch(waveBody.frequency);
    vibrato.patch(waveString1.frequency);
    vibrato.patch(waveString2.frequency);

    mix = new Summer();
    waveBody.patch(mix);
    waveString1.patch(mix);
    waveString2.patch(mix);

    // フィルタ（つまみ：明るさ gBright・共振 gReso）。固定（スイープ無し）
    filter = new MoogFilter(gBright, gReso, MoogFilter.Type.LP);
    mix.patch(filter);

    // 弓ノイズ・エンベロープ（弾き方ごと）。擦れ音の量は つまみ gBowNoise
    float attackNoiseAmp, frictionNoiseAmp;
    if (art == 1) {            // 上げ弓（軽い）
      attackNoiseAmp   = amplitude * 0.30;
      frictionNoiseAmp = amplitude * gBowNoise;
      attackEnv   = new ADSR(1.0, 0.016 * lenF, 0.04, 0.0,  0.07 * lenF);
      frictionEnv = new ADSR(1.0, 0.025,        0.05, 0.85, 0.09 * lenF);
      ampEnv      = new ADSR(1.0, 0.022 * lenF, 0.06, 0.82, 0.10 * lenF);
    } else if (art == 2) {     // デタシェ
      attackNoiseAmp   = amplitude * 0.34;
      frictionNoiseAmp = amplitude * gBowNoise;
      attackEnv   = new ADSR(1.0, 0.010 * lenF, 0.03, 0.0,  0.05 * lenF);
      frictionEnv = new ADSR(1.0, 0.015,        0.03, 0.90, 0.06 * lenF);
      ampEnv      = new ADSR(1.0, 0.014 * lenF, 0.04, 0.80, 0.06 * lenF);
    } else {                   // 下げ弓（しっかり）
      attackNoiseAmp   = amplitude * 0.38;
      frictionNoiseAmp = amplitude * gBowNoise;
      attackEnv   = new ADSR(1.0, 0.020 * lenF, 0.05, 0.0,  0.10 * lenF);
      frictionEnv = new ADSR(1.0, 0.045,        0.07, 0.88, 0.14 * lenF);
      ampEnv      = new ADSR(1.0, 0.040 * lenF, 0.08, 0.85, 0.16 * lenF);
    }

    bowAttack   = new Noise(attackNoiseAmp,   Noise.Tint.WHITE);
    bowFriction = new Noise(frictionNoiseAmp, Noise.Tint.PINK); // 弓の擦れ音

    bowAttack.patch(attackEnv);
    attackEnv.patch(mix);
    bowFriction.patch(frictionEnv);
    frictionEnv.patch(mix);
    filter.patch(ampEnv);
  }

  void noteOn(float duration) {
    ampEnv.noteOn();
    attackEnv.noteOn();
    frictionEnv.noteOn();
    pitchSlide.activate();
    ampEnv.patch(out);
  }

  void noteOff() {
    ampEnv.noteOff();
    attackEnv.noteOff();
    frictionEnv.noteOff();
    ampEnv.unpatchAfterRelease(out);
  }
}
\end{lstlisting}